\section{Criterios diferenciales de lisitud formal y de regularidad}
\label{section:0.22-es}

Los principales resultados de este parágrafo son los siguientes:
\begin{enumerate}
  \item[a)] Los criterios (22.2.2) y (22.5.8), que dan condiciones
  necesarias y suficientes para que un anillo local noetheriano completo $A$
  que contiene un cuerpo $k$ sea formalmente liso sobre $k$; el criterio
  (22.5.8) se enuncia sin hacer intervenir nociones diferenciales y completa
  (19.6.6); el enunciado (22.2.2), de naturaleza algo más técnica, permite dar
  una \emph{clasificación}, para $k$ fijado, de los anillos locales $A$
  considerados (22.2.5): vienen dados por su dimensión $n$ y por su cuerpo
  residual $K$, con la única condición
  $n\geqslant\operatorname{rg}_K\Upsilon_{K/k}$.

  \item[b)] El teorema (22.3.3) y el corolario (22.7.6) del criterio
  jacobiano de Nagata (22.7.3), que dan condiciones que permiten afirmar que
  una \emph{localización} de un anillo local noetheriano completo $A$ que
  contiene un cuerpo $k$ es geométricamente regular sobre $k$; desempeñarán
  un papel importante en el estudio de las propiedades «~finas~» de los
  anillos locales realizado en el cap.~IV, §~7. Por el momento no se conoce
  ninguna demostración de (22.7.6) independiente del criterio jacobiano de
  Nagata.

  \item[c)] El criterio jacobiano de Zariski (22.6.7) y sus variantes, que
  son consecuencias fáciles de (20.7.8).
\end{enumerate}

\oldpage[0\textsubscript{IV-1}]{183}

\subsection{Levantamiento de la lisitud formal}
\label{subsection:0.22.1-es}

\begin{theorem}[22.1.1]
\label{0.22.1.1-es}
Sean $s:\Lambda\to A$, $u:A\to B$ dos homomorfismos de anillos,
$\mathfrak J$ un ideal de $B$ y $C$ el anillo cociente $B/\mathfrak J$. Se
supone que $C$ es una $\Lambda$-álgebra formalmente lisa para la topología
discreta.
\begin{enumerate}
  \item[(i)] Supongamos que se cumplen las condiciones siguientes:
  \begin{enumerate}
    \item[$\alpha$)] $\mathfrak J/\mathfrak J^2$ es un $C$-módulo proyectivo.
    \item[$\beta$)] El homomorfismo canónico
    $\mathbf S_C^*(\mathfrak J/\mathfrak J^2)\to
    \operatorname{gr}_{\mathfrak J}^*(B)$ (19.5.2) es biyectivo.
    \item[$\gamma$)] El homomorfismo característico
    $\chi:\Upsilon_{C/B/\Lambda}\to\mathfrak J/\mathfrak J^2$ (20.6.20)
    es inyectivo.
    \item[$\delta$)] El $C$-módulo
    $\Omega^1_{B/A}\otimes_BC$ es proyectivo.
  \end{enumerate}
  Entonces $B$, provisto de la topología $\mathfrak J$-preádica, es una
  $\Lambda$-álgebra formalmente lisa.

  \item[(ii)] Recíprocamente, supongamos que $B$ sea una $\Lambda$-álgebra
  formalmente lisa para la topología $\mathfrak J$-preádica y que $A$ sea una
  $\Lambda$-álgebra formalmente lisa para la topología discreta. Entonces se
  cumplen las condiciones $\alpha$), $\beta$), $\gamma$), $\delta$) de (i).
\end{enumerate}
\end{theorem}

\begin{proof}
En virtud de la sucesión exacta (20.6.22.1) (que es aplicable, puesto que
$B/\mathfrak J^2$ es una extensión $\Lambda$-trivial de $C$ por
$\mathfrak J/\mathfrak J^2$, al suponerse que $C$ es una $\Lambda$-álgebra
formalmente lisa (19.4.4)), la condición $\gamma$) equivale a
$\Upsilon^C_{B/A/\Lambda}=0$, o también a la siguiente:
\begin{enumerate}
  \item[$\gamma'$)] El homomorfismo
  \[
    u_{B/A/\Lambda}\otimes 1_C:
    \Omega^1_{A/\Lambda}\otimes_A C
    \longrightarrow \Omega^1_{B/\Lambda}\otimes_B C
  \]
  es inyectivo.
\end{enumerate}

(i) Las condiciones $\alpha$) y $\beta$) implican que $B$ es una
$\Lambda$-álgebra formalmente lisa para la topología
$\mathfrak J$-preádica (19.5.4). Por tanto, equivale decir que $B$ es (para la
topología $\mathfrak J$-preádica) una $A$-álgebra formalmente lisa o una
$A$-álgebra formalmente lisa \emph{relativamente a $\Lambda$}. Para ver que
$B$ es una $A$-álgebra formalmente lisa basta, en virtud de (20.7.2), probar
que el homomorfismo continuo de $B$-módulos topológicos
\[
  u_{B/A/\Lambda}:\Omega^1_{A/\Lambda}\otimes_A B
  \longrightarrow\Omega^1_{B/\Lambda}
\]
es \emph{formalmente invertible por la izquierda}. Ahora bien, puesto que $B$
es una $\Lambda$-álgebra formalmente lisa, el $B$-módulo topológico
$\Omega^1_{B/\Lambda}$ es formalmente proyectivo (20.4.9) y su topología se
deduce de la de $B$ (20.4.5). En virtud de (19.1.9), para que
$u_{B/A/\Lambda}$ sea formalmente invertible por la izquierda basta que
$u_{B/A/\Lambda}\otimes 1_C$ sea invertible por la izquierda. Pero, en virtud
de $\gamma'$), esta última aplicación es inyectiva, por lo que se tiene la
sucesión exacta (20.6.14.7)
\[
  0\longrightarrow\Omega^1_{A/\Lambda}\otimes_A C
  \longrightarrow\Omega^1_{B/\Lambda}\otimes_B C
  \longrightarrow\Omega^1_{B/A}\otimes_B C\longrightarrow0.
\]
Finalmente, en virtud de $\delta$), el $C$-módulo
$\Omega^1_{B/A}\otimes_BC$ es proyectivo, luego la sucesión exacta anterior
se escinde, lo que termina de probar (i).

\oldpage[0\textsubscript{IV-1}]{184}

(ii) Si $B$ es una $\Lambda$-álgebra formalmente lisa para la topología
$\mathfrak J$-preádica y $A$ una $\Lambda$-álgebra formalmente lisa para la
topología discreta, $B$ es una $A$-álgebra formalmente lisa para la topología
$\mathfrak J$-preádica (19.3.5, (ii)); por tanto, las condiciones $\alpha$) y
$\beta$) resultan de (19.5.6). Además, (20.7.2) muestra que
$u_{B/A/\Lambda}$ es formalmente invertible por la izquierda, luego
$u_{B/A/\Lambda}\otimes 1_C$ es invertible por la izquierda (19.1.7) y, \emph{a
fortiori}, inyectivo. Finalmente, $\Omega^1_{B/A}$ es un $B$-módulo
formalmente proyectivo (20.4.9) y, por consiguiente,
$\Omega^1_{B/A}\otimes_BC$ es un $C$-módulo proyectivo (19.2.4).
\end{proof}

\begin{corollary}[22.1.2]
\label{0.22.1.2-es}
Sean $s:\Lambda\to A$, $u:A\to B$ dos homomorfismos de anillos,
$\mathfrak m$ un ideal maximal de $B$ y $K=B/\mathfrak m$ el cuerpo cociente.
Se supone que las $\Lambda$-álgebras $A$ y $K$ son formalmente lisas para la
topología discreta. Entonces las condiciones siguientes son equivalentes:
\begin{enumerate}
  \item[a)] $B$ es una $A$-álgebra formalmente lisa para la topología
  $\mathfrak m$-preádica.
  \item[b)] El homomorfismo canónico
  \begin{equation}
  \tag{22.1.2.1}
  \label{0.22.1.2.1-es}
    \mathbf S_K^*(\mathfrak m/\mathfrak m^2)
    \longrightarrow\operatorname{gr}_{\mathfrak m}^*(B)
  \end{equation}
  es biyectivo, y el homomorfismo característico
  \begin{equation}
  \tag{22.1.2.2}
  \label{0.22.1.2.2-es}
    \chi:\Upsilon_{K/B/\Lambda}\longrightarrow\mathfrak m/\mathfrak m^2
  \end{equation}
  es inyectivo.
\end{enumerate}
\end{corollary}

\begin{proof}
Esto resulta inmediatamente de (22.1.1), pues las condiciones $\alpha$) y
$\delta$) se cumplen aquí trivialmente, ya que $K$ es un cuerpo.
\end{proof}

\begin{remark}[22.1.3]
\label{0.22.1.3-es}
Supongamos que se cumplen las hipótesis de (22.1.2) y, además, que
$\mathfrak m/\mathfrak m^2$ es un $K$-espacio vectorial de rango \emph{finito}.
Entonces, para que $B$ sea una $A$-álgebra formalmente lisa para la topología
$\mathfrak m$-preádica, basta que se satisfagan las condiciones siguientes:

\begin{env}[22.1.3.1]
\label{0.22.1.3.1-es}
Dados un álgebra de polinomios $E=K[T_1,\ldots,T_n]$, el ideal maximal
$\mathfrak n$ de $E$ engendrado por los $T_i$ $(1\leqslant i\leqslant n)$ y
un ideal $\mathfrak b\supset\mathfrak n^k$ de $E$, todo
$\Lambda$-homomorfismo $v:B\to E/\mathfrak b$ que, al pasar a los cocientes,
da la identidad $K\to K$, se factoriza como
\[
  B\xrightarrow{w}E/\mathfrak n^k\longrightarrow E/\mathfrak b,
\]
donde $w$ es un $\Lambda$-homomorfismo.
\end{env}

\begin{env}[22.1.3.2]
\label{0.22.1.3.2-es}
Si $F$ es la extensión trivial $D_K(K)$ («~álgebra de los números duales
sobre $K$~»), $\rho:A\to F$ un homomorfismo de anillos tal que el compuesto
$A\xrightarrow{\rho}F\to K$ sea el homomorfismo canónico, entonces todo
$\Lambda$-homomorfismo $v:B\to K$ que, al pasar a los cocientes, da la
identidad $K\to K$, se factoriza como
\[
  B\xrightarrow{w}F\longrightarrow K,
\]
donde $w$ es un $\Lambda$-homomorfismo (para la estructura de $A$-álgebra
sobre $F$ definida por $\rho$).
\end{env}

En efecto, se ha visto (19.5.5) que la condición (22.1.3.1) implica que el
homomorfismo canónico (22.1.2.1) es biyectivo. En virtud de (22.1.2), basta
entonces ver que el homomorfismo canónico
\[
  u_{B/A/\Lambda}\otimes1_K:
  \Omega^1_{A/\Lambda}\otimes_AK
  \longrightarrow\Omega^1_{B/\Lambda}\otimes_BK
\]
es inyectivo, o también (puesto que se trata de espacios vectoriales sobre
$K$) que, para todo $K$-espacio vectorial $L$, el homomorfismo canónico
$\operatorname{D\acute{e}r}_\Lambda(B,L)\to\operatorname{D\acute{e}r}_\Lambda(A,L)$ es
\emph{sobreyectivo} (en virtud de (20.4.8)); es claro que, para ello, es
necesario y suficiente que el homomorfismo canónico
$\operatorname{D\acute{e}r}_\Lambda(B,K)\to\operatorname{D\acute{e}r}_\Lambda(A,K)$ sea
sobreyectivo, de donde resulta nuestra afirmación en virtud de (20.1.5).

Se observará que las dos condiciones (22.1.3.1) y (22.1.3.2) son consecuencia
de la siguiente:

\begin{env}[22.1.3.3]
\label{0.22.1.3.3-es}
Para toda $\Lambda$-álgebra local artiniana $E$ de cuerpo residual igual a
$K$, y todo ideal nilpotente $\mathfrak r$ de $E$, todo
$\Lambda$-homomorfismo $B\to E/\mathfrak r$ que, al pasar a los cocientes,
da la identidad $K\to K$, se factoriza como
$B\xrightarrow{u}E\to E/\mathfrak r$, donde $u$ es un
$\Lambda$-homomorfismo.
\end{env}
\end{remark}

\oldpage[0\textsubscript{IV-1}]{185}

Este resultado se generaliza como sigue:

\begin{proposition}[22.1.4]
\label{0.22.1.4-es}
Sea $\varphi:A\to B$ un homomorfismo local de anillos locales noetherianos,
$k$ (resp. $K$) el cuerpo residual de $A$ (resp. de $B$). Para que $B$ sea una
$A$-álgebra formalmente lisa (para las topologías preádicas), es necesario y
suficiente que, para todo anillo local artiniano $C$, de cuerpo residual igual
a $K$, todo homomorfismo local $A\to C$ (que hace de $C$ una $A$-álgebra
topológica), y todo ideal nilpotente $\mathfrak J$ de $C$, todo
$A$-homomorfismo local $B\to C/\mathfrak J$ tal que el homomorfismo
$K\to K$ obtenido al pasar a los cocientes sea la identidad, se factorice como
\[
  B\xrightarrow{f}C\longrightarrow C/\mathfrak J,
\]
donde $f$ es un $A$-homomorfismo local.
\end{proposition}

\begin{proof}
Como la condición es necesaria por definición (19.3.1), mostremos que es
suficiente. Puesto que los $\widehat A$-homomorfismos locales de $\widehat B$
en una $\widehat A$-álgebra local discreta proceden, por extensión, de los
$A$-homomorfismos locales de $B$ en esta álgebra, podemos limitarnos al caso
en que $A$ y $B$ sean \emph{completos}, en virtud de (19.3.6). Cuando $A$ es
un \emph{cuerpo} $k$, la proposición resulta de (22.1.3.3), tomando para
$\Lambda$ el subcuerpo primo de $k$, y de (19.6.1).

En el caso general, pongamos $B_0=B\otimes_Ak=B/\mathfrak mB$
($\mathfrak m$ ideal maximal de $A$), que es completo. Si $C_0$ es una
$k$-álgebra local artiniana de cuerpo residual igual a $K$, $\mathfrak J_0$
un ideal nilpotente de $C_0$, todo $k$-homomorfismo local
$g_0:B_0\to C_0/\mathfrak J_0$ que da la identidad sobre $K$ al pasar a los
cocientes da, por composición, un $A$-homomorfismo local
\[
  g:B\longrightarrow B_0\xrightarrow{g_0}C_0/\mathfrak J_0
\]
que, por hipótesis, se factoriza por tanto como
$B\xrightarrow{f}C_0\to C_0/\mathfrak J_0$, donde $f$ es un
$A$-homomorfismo local; pero, como $C_0$ es una $k$-álgebra, $f$ se factoriza
como $B\to B_0\xrightarrow{f_0}C_0$, donde $f_0$ es un
$k$-homomorfismo local. Se ve así que las hipótesis del enunciado siguen
cumpliéndose cuando se sustituyen $B_0$ y $k$ por $B$ y $A$, respectivamente;
por tanto, $B_0$ es una $k$-álgebra \emph{formalmente lisa} por lo que se acaba
de ver.

Aplicando (19.7.2) y (19.7.1), se ve que existen un anillo local noetheriano
completo $B'$, un homomorfismo local $A\to B'$ que hace de $B'$ un $A$-módulo
\emph{plano} y una $A$-álgebra \emph{formalmente lisa}, y un
$k$-isomorfismo
\[
  u'_0:B'_0=B'\otimes_Ak\xrightarrow{\sim}B_0.
\]
Sea $v'_0$ el isomorfismo recíproco de $u'_0$; nos proponemos mostrar que
existe un $A$-homomorfismo local $v:B\to B'$ tal que $v'_0$ proceda de $v$ al
pasar a los cocientes. Para ello, designemos por $\mathfrak n$ y
$\mathfrak n'$ los ideales maximales respectivos de $B$ y $B'$ y, para todo
$j$, sea
\[
  w_j:B/(\mathfrak n^{j+1}+\mathfrak mB)
  \longrightarrow B'/(\mathfrak n'^{j+1}+\mathfrak mB')
\]
el homomorfismo deducido de $v'_0$ al pasar a los cocientes; bastará construir,
para todo $j$, un $A$-homomorfismo local
$v_j:B/\mathfrak n^{j+1}\to B'/\mathfrak n'^{j+1}$ tal que los diagramas
\[
\xymatrix@C=4em@R=3em{
  B/\mathfrak n^{j+1} \ar[r]^{v_j} \ar[d] &
  B'/\mathfrak n'^{j+1} \ar[d] \\
  B/\mathfrak n^j \ar[r]_{v_{j-1}} & B'/\mathfrak n'^j
}
\qquad
\xymatrix@C=4em@R=3em{
  B/\mathfrak n^{j+1} \ar[r]^{v_j} \ar[d] &
  B'/\mathfrak n'^{j+1} \ar[d] \\
  B/(\mathfrak n^{j+1}+\mathfrak mB) \ar[r]_{w_j} &
  B'/(\mathfrak n'^{j+1}+\mathfrak mB')
}
\]
sean conmutativos; como $B$ y $B'$ son completos, el homomorfismo
$v=\varprojlim v_j$ responderá a la cuestión.

\oldpage[0\textsubscript{IV-1}]{186}

Ahora bien, la construcción por recurrencia de $v_j$ resulta de la hipótesis
sobre $B$ y del mismo razonamiento que en (19.3.10.1), utilizando el lema
(19.3.10.2) y el hecho de que
\[
  \mathfrak n'^j+(\mathfrak n'^{j+1}+\mathfrak mB')
  =\mathfrak n'^j+\mathfrak mB'.
\]
Resulta entonces de (19.7.1.4) que $v$ es un $A$-\emph{isomorfismo}, pues
$B'$ es un $A$-módulo plano y $B$ es completo para la topología
$\mathfrak m$-preádica, pues los ideales $\mathfrak m^iB$ son cerrados en $B$
para la topología $\mathfrak n$-ádica ($\mathbf 0_{\mathrm I}$, 7.3.5).
Q.E.D.
\end{proof}

\subsection{Caracterización diferencial de las álgebras locales formalmente
lisas sobre un cuerpo}
\label{subsection:0.22.2-es}

\begin{env}[22.2.1]
\label{0.22.2.1-es}
Sean $k$ un cuerpo, $P$ su subcuerpo primo, $A$ una $k$-álgebra que es un
anillo \emph{local}, $\mathfrak m$ su ideal maximal y $K=A/\mathfrak m$ su
cuerpo residual. Como $K$ es separable sobre $P$, $K$ es una $P$-álgebra
formalmente lisa para la topología discreta (19.6.1), y por consiguiente la
$P$-extensión $A/\mathfrak m^2$ de $K$ por $\mathfrak m/\mathfrak m^2$ es
$P$-\emph{trivial}, y el homomorfismo característico
\begin{equation}
\tag{22.2.1.1}
\label{0.22.2.1.1-es}
  \chi_{A/k}:\Upsilon_{K/k}\longrightarrow\mathfrak m/\mathfrak m^2
\end{equation}
queda así definido (20.6.23).
\end{env}

\begin{theorem}[22.2.2]
\label{0.22.2.2-es}
Bajo las condiciones de (22.2.1), para que $A$ sea una $k$-álgebra
formalmente lisa para la topología $\mathfrak m$-preádica, es necesario y
suficiente que se satisfagan las condiciones siguientes:
\begin{enumerate}
  \item[(i)] El homomorfismo canónico
  $\mathbf S_K^*(\mathfrak m/\mathfrak m^2)\to
  \operatorname{gr}_{\mathfrak m}^*(A)$ es biyectivo.
  \item[(ii)] El homomorfismo característico
  $\chi_{A/k}:\Upsilon_{K/k}\to\mathfrak m/\mathfrak m^2$ es inyectivo.
\end{enumerate}
\end{theorem}

\begin{proof}
Es un caso particular de (22.1.2), donde se sustituye $\Lambda$ por $P$, $A$
por $k$ y $B$ por $A$; en efecto, como $K$ y $k$ son separables sobre $P$,
son $P$-álgebras formalmente lisas (19.6.1).
\end{proof}

\begin{remark}[22.2.3]
\label{0.22.2.3-es}
La condición (ii) de (22.2.2) puede sustituirse por cualquiera de las
siguientes:
\begin{enumerate}
  \item[a)] El homomorfismo canónico
  $\Omega^1_k\otimes_kK\to\Omega^1_A\otimes_AK$ es inyectivo.
  \item[b)] Para todo $n$, el homomorfismo canónico
  \[
    \Omega^1_k\otimes_k(A/\mathfrak m^{n+1})
    \longrightarrow
    \Omega^1_A\otimes_A(A/\mathfrak m^{n+1})
  \]
  es invertible por la izquierda.
  \item[c)] El homomorfismo canónico
  $\Omega^1_k\widehat\otimes_kA\to\widehat\Omega^1_A$ es invertible por la
  izquierda.
  \item[d)] $A$ es una $k$-álgebra formalmente lisa (para la topología
  $\mathfrak m$-preádica) \emph{relativamente al cuerpo primo}.
\end{enumerate}
En efecto, ya se ha observado en (22.1.1) que la inyectividad de $\chi_{A/k}$
equivale a $a)$, y que la conjunción de $a)$ y de la condición (i) de (22.2.2)
equivale a la de $d)$ y (i). Finalmente, también se ha visto en (22.1.1) que
$b)$ y $d)$ son equivalentes, y la equivalencia de $b)$ y $c)$ resulta de
(19.1.8).
\end{remark}

\begin{env}[22.2.4]
\label{0.22.2.4-es}
La principal aplicación de (22.2.2) se refiere a las $k$-álgebras locales
noetherianas y \emph{completas} (caso en que $\mathfrak m/\mathfrak m^2$ es un
$K$-espacio vectorial de rango \emph{finito}). Precisamente, sean $k$ un
cuerpo, $K$ una extensión de $k$ y $V$ un $K$-espacio vectorial de

\oldpage[0\textsubscript{IV-1}]{187}

rango finito; consideramos las ternas $(A,\alpha,\beta)$, donde $A$ es una
$k$-álgebra local noetheriana completa \emph{formalmente lisa} de ideal
maximal $\mathfrak m$, $\alpha:A/\mathfrak m\xrightarrow{\sim}K$ un
isomorfismo de $k$-álgebras y
$\beta:\mathfrak m/\mathfrak m^2\xrightarrow{\sim}V$ un diisomorfismo de
espacios vectoriales (para el isomorfismo $\alpha$). Dadas dos de tales ternas
$(A,\alpha,\beta)$, $(A',\alpha',\beta')$, una \emph{equivalencia} de la
primera sobre la segunda es un $k$-isomorfismo $u:A\xrightarrow{\sim}A'$ tal
que (si $\mathfrak m'$ es el ideal maximal de $A'$) los diagramas
\[
\xymatrix@C=4em@R=3em{
  A/\mathfrak m \ar[r]^{\operatorname{gr}^0(u)} \ar[d]_{\alpha} &
  A'/\mathfrak m' \ar[d]^{\alpha'} \\
  K \ar[r]_{1_K} & K
}
\qquad
\xymatrix@C=4em@R=3em{
  \mathfrak m/\mathfrak m^2 \ar[r]^{\operatorname{gr}^1(u)}
    \ar[d]_{\beta} &
  \mathfrak m'/\mathfrak m'^2 \ar[d]^{\beta'} \\
  V \ar[r]_{1_V} & V
}
\]
sean conmutativos. Es inmediato que las \emph{clases de equivalencia} de las
ternas $(A,\alpha,\beta)$ forman un \emph{conjunto} que denotaremos por
$\operatorname{FL}(K/k,V)$. Observemos ahora que a toda terna
$(A,\alpha,\beta)$ se asocia el $K$-homomorfismo compuesto de espacios
vectoriales
\[
  \Upsilon_{K/k}\xrightarrow{\chi_{A/k}}
  \mathfrak m/\mathfrak m^2\xrightarrow{\beta}V
\]
(donde $\chi_{A/k}$ se considera como un dihomomorfismo relativo a
$\alpha^{-1}$); la definición anterior demuestra (por (20.6.22.2)) que este
$K$-homomorfismo \emph{solo depende de la clase de equivalencia} de
$(A,\alpha,\beta)$ en $\operatorname{FL}(K/k,V)$; en otras palabras, se ha
definido una aplicación canónica
\begin{equation}
\tag{22.2.4.1}
\label{0.22.2.4.1-es}
  \operatorname{FL}(K/k,V)\longrightarrow
  \operatorname{Hom}_K(\Upsilon_{K/k},V).
\end{equation}
\end{env}

\begin{proposition}[22.2.5]
\label{0.22.2.5-es}
Sean $k$ un cuerpo, $K$ una extensión de $k$ y $V$ un $K$-espacio vectorial
de rango finito. La aplicación canónica (22.2.4.1) es una biyección de
$\operatorname{FL}(K/k,V)$ sobre el conjunto de los $K$-homomorfismos
inyectivos de $\Upsilon_{K/k}$ en $V$.
\end{proposition}

\begin{proof}
(i) Para mostrar que (22.2.4.1) es inyectiva, hay que probar lo siguiente:
sean $A$, $A'$ dos $k$-álgebras locales noetherianas completas formalmente
lisas, $\mathfrak m$, $\mathfrak m'$ sus ideales maximales respectivos,
$\alpha:A/\mathfrak m\xrightarrow{\sim}K$,
$\alpha':A'/\mathfrak m'\xrightarrow{\sim}K$ dos $k$-isomorfismos;
supongamos dados dos $k$-isomorfismos
\[
  v^0:A/\mathfrak m\xrightarrow{\sim}A'/\mathfrak m',\qquad
  v^1:\mathfrak m/\mathfrak m^2\xrightarrow{\sim}
      \mathfrak m'/\mathfrak m'^2
\]
tales que los diagramas
\begin{equation}
\tag{22.2.5.1}
\label{0.22.2.5.1-es}
\xymatrix@C=4em@R=3em{
  A/\mathfrak m \ar[r]^{v^0} \ar[dr]_{\alpha} &
  A'/\mathfrak m' \ar[dl]^{\alpha'} \\
  & K
}
\qquad
\xymatrix@C=4em@R=3em{
  \mathfrak m/\mathfrak m^2 \ar[r]^{v^1} &
  \mathfrak m'/\mathfrak m'^2 \\
  \Upsilon_{K/k} \ar[u]^{\chi_A} \ar[ur]_{\chi_{A'}} &
}
\end{equation}
sean conmutativos. Se trata de probar que existe un $k$-isomorfismo
$u:A\xrightarrow{\sim}A'$ tal que $\operatorname{gr}^0(u)=v^0$ y
$\operatorname{gr}^1(u)=v^1$. Observemos primero que, puesto que $K$ es un
cuerpo, el homomorfismo canónico (20.6.7.1) es \emph{biyectivo} (20.6.11),
luego ya existe un $k$-isomorfismo
$v:A/\mathfrak m^2\xrightarrow{\sim}A'/\mathfrak m'^2$ tal que
$\operatorname{gr}^0(v)=v^0$ y $\operatorname{gr}^1(v)=v^1$. Observemos ahora
que, puesto que $A'$ es metrizable y completo y $A$ una $k$-álgebra
formalmente lisa, el homomorfismo compuesto
$A\to A/\mathfrak m^2\xrightarrow{v}A'/\mathfrak m'^2$ se factoriza como
$A\xrightarrow{u}A'\to A'/\mathfrak m'^2$, donde $u$ es

\oldpage[0\textsubscript{IV-1}]{188}

un $k$-homomorfismo (19.3.11), y se tiene ya, por tanto,
$\operatorname{gr}^0(u)=v^0$ y $\operatorname{gr}^1(u)=v^1$. Ahora bien, se
tiene el diagrama conmutativo
\[
\xymatrix@C=5em@R=3em{
  \mathbf S_K^*(\mathfrak m/\mathfrak m^2) \ar[r] \ar[d]_{w} &
  \operatorname{gr}_{\mathfrak m}^*(A) \ar[d]^{\operatorname{gr}^*(u)} \\
  \mathbf S_K^*(\mathfrak m'/\mathfrak m'^2) \ar[r] &
  \operatorname{gr}_{\mathfrak m'}^*(A')
}
\]
donde las flechas horizontales son los homomorfismos canónicos y $w$ procede
de $v^0$ y $v^1$; como $v^0$ y $v^1$ son biyectivos, también lo es $w$, y,
puesto que $A$ y $A'$ son $k$-álgebras formalmente lisas, se deduce de
(22.2.2, (i)) que $\operatorname{gr}(u)$ es biyectivo; pero, como $A$ y $A'$
son separados y completos, esto implica que el propio $u$ es
\emph{biyectivo} (Bourbaki, \emph{Alg. comm.}, cap.~III, §~2,
n\textsuperscript{o}~8, cor.~3 del th.~1).

(ii) Queda por mostrar que la imagen de (22.2.4.1) es el conjunto de los
homomorfismos \emph{inyectivos} de $\Upsilon_{K/k}$ en $V$; ya se sabe por
(22.2.2, (ii)) que está contenida en este conjunto; bastará entonces probar el
lema siguiente.
\end{proof}

\begin{lemma}[22.2.5.2]
\label{0.22.2.5.2-es}
Para todo $K$-homomorfismo $h:\Upsilon_{K/k}\to V$, existe una $k$-álgebra
$A$ que es un anillo local noetheriano regular y completo, de ideal maximal
$\mathfrak m$, de cuerpo residual $K$, y un $K$-isomorfismo
$\beta:\mathfrak m/\mathfrak m^2\xrightarrow{\sim}V$ tales que
$h=\beta\circ\chi_{A/k}$.
\end{lemma}

\begin{proof}
En efecto, si se prueba este lema, el hecho de que $A$ sea regular implica que
se satisface la condición (i) de (22.2.2) (17.1.1); si además $h$ es
inyectivo, (22.2.2) mostrará que $A$ es una $k$-álgebra formalmente lisa cuya
clase tendrá por imagen $h$ (en cambio, si $h$ no es inyectivo, la $k$-álgebra
$A$ cuya existencia prueba (22.2.5.2) no es formalmente lisa).

Para probar (22.2.5.2), sean $B$ la $K$-álgebra $\mathbf S_K^*(V)$,
$\mathfrak n$ el ideal de aumentación de $B$ y $A=\widehat B$ la $K$-álgebra
completada de $B$ para la topología $\mathfrak n$-preádica (de modo que $A$ es
isomorfa al álgebra de series formales $K[[T_1,\ldots,T_n]]$ si
$n=\operatorname{rg}_KV$). En virtud de (20.6.11), existe sobre
$A/\mathfrak m^2$ (donde $\mathfrak m=\mathfrak nA$) una estructura de
$k$-extensión de $K$ por $V$ (en general \emph{distinta} de la estructura
$k$-trivial definida por la inyección canónica $k\to K\to A$) tal que $h$ sea
el homomorfismo característico de esta extensión. Si
$f:k\to A/\mathfrak m^2$ es el homomorfismo estructural, el hecho de que $k$
sea separable sobre el cuerpo primo permite (19.6.1) factorizar $f$ como
$k\xrightarrow{g}A\to A/\mathfrak m^2$, y la $k$-álgebra $A$ así definida
responde a la cuestión (20.6.22.2).
\end{proof}

\begin{theorem}[22.2.6]
\label{0.22.2.6-es}
Sean $k$ un cuerpo y $K$ una extensión de $k$. Para que exista una $k$-álgebra
local noetheriana $A$, de ideal maximal $\mathfrak m$, tal que
$A/\mathfrak m=K$ y que $A$ sea formalmente lisa (para su topología
$\mathfrak m$-preádica), es necesario y suficiente que $\Upsilon_{K/k}$ sea
un $K$-espacio vectorial de rango finito. La $k$-álgebra $A$ queda entonces
determinada (salvo un $k$-isomorfismo que, al pasar a los cocientes, da la
identidad $K\to K$) por su dimensión, que está sometida a la única condición
de ser $\geqslant\operatorname{rg}_K\Upsilon_{K/k}$.
\end{theorem}

\begin{proof}
Como $\mathfrak m/\mathfrak m^2$ es de rango finito sobre $K$, la necesidad
de la condición resulta de (22.2.2, (ii)); recíprocamente, (22.2.5) muestra que
la condición es suficiente y que se puede tomar como
$\mathfrak m/\mathfrak m^2$ un $K$-espacio vectorial de rango cualquiera
$\geqslant\operatorname{rg}_K(\Upsilon_{K/k})$.
\end{proof}

\oldpage[0\textsubscript{IV-1}]{189}

En particular, el homomorfismo identidad de $\Upsilon_{K/k}$ asocia
\emph{canónicamente} a $K$ una $k$-álgebra noetheriana completa, formalmente
lisa, para la cual $\Upsilon_{K/k}$ es isomorfo a
$\mathfrak m/\mathfrak m^2$.

\begin{remarks}[22.2.7]
\label{0.22.2.7-es}
\begin{enumerate}
  \item[(i)] Sea $\Lambda$ un anillo local noetheriano, de cuerpo residual
  $k$; resulta de (19.7.1) y (19.7.2) que la determinación de las
  $\Lambda$-álgebras $A$ que son anillos locales noetherianos completos y que
  son formalmente lisas es \emph{equivalente} al mismo problema en el que
  $\Lambda$ se sustituye por $k$, y queda por tanto, en principio, enteramente
  resuelta por (22.2.5), que reduce la cuestión a la estructura de los
  \emph{cuerpos residuales} de las $k$-álgebras buscadas, como extensiones de
  $k$.

  \item[(ii)] El hecho de que $\Upsilon_{K/k}$ sea de rango finito sobre $K$
  \emph{no implica} que $K$ sea «~de multiplicidad radicial finita~» sobre
  $k$ (cf. (19.6.6)). Por ejemplo, sean $k_0$ un cuerpo perfecto de
  característica $p>0$, $k=k_0(T)$ el cuerpo de las fracciones racionales en
  una indeterminada sobre $k_0$ y $K=k^{p^{-\infty}}$ la menor extensión
  perfecta de $k$. Se tiene entonces $\Omega^1_K=\Omega^1_{K/k}=0$ (21.4.4),
  luego, por definición (20.6.1.1),
  $\Upsilon_{K/k}=\Omega^1_k\otimes_kK$; y como $k^p=k_0(T^p)$, $T$ es una
  $p$-base de $k$ sobre $k^p$, luego (21.4.5) $\Omega^1_k$ es de rango $1$
  sobre $k$ y, por consiguiente, $\Upsilon_{K/k}$ de rango $1$ sobre $K$.
  Pero, para todo $n$, el cuerpo residual del anillo local
  $K\otimes_kk^{p^{-n}}$ es isomorfo a $K$ y, por tanto, no es separable sobre
  $k^{p^{-n}}$.
\end{enumerate}
\end{remarks}

El teorema (22.2.2) se generaliza como sigue:

\begin{proposition}[22.2.8]
\label{0.22.2.8-es}
Bajo las condiciones de (22.2.1), supongamos que la característica $p$ de $k$
sea $>0$; sea $(k_\alpha)$ una familia filtrante decreciente de subcuerpos de
$k$ tal que $\bigcap_\alpha k_\alpha(k^p)=k^p$. Entonces la condición (ii) de
(22.2.2) puede sustituirse por cualquiera de las siguientes:
\begin{enumerate}
  \item[a)] Existe $\alpha_0$ tal que el homomorfismo canónico
  \[
    \Omega^1_{k/k_\alpha}\otimes_kK
    \longrightarrow\Omega^1_{A/k_\alpha}\otimes_AK
  \]
  sea inyectivo para todo $\alpha\geqslant\alpha_0$.
  \item[b)] Existe $\alpha_0$ tal que $A$ sea una $k$-álgebra formalmente
  lisa (para la topología $\mathfrak m$-preádica) relativamente a $k_\alpha$,
  para todo $\alpha\geqslant\alpha_0$.
  \item[c)] Para todo $\alpha$ y todo entero $n\geqslant0$, el homomorfismo
  canónico
  \[
    \Omega^1_{k/k_\alpha}\otimes_k(A/\mathfrak m^{n+1})
    \longrightarrow
    \Omega^1_{A/k_\alpha}\otimes_A(A/\mathfrak m^{n+1})
  \]
  es invertible por la izquierda.
\end{enumerate}
\end{proposition}

\begin{proof}
En efecto, si $A$ es una $k$-álgebra formalmente lisa, lo es \emph{a
fortiori} relativamente a todo subcuerpo de $k$, lo que implica $c)$ en virtud
de (20.7.2) y (19.1.7); es claro que $c)$ implica $b)$, en virtud de (20.7.2),
y que $b)$ implica $a)$. Finalmente, teniendo en cuenta (22.2.3), queda por
probar que $a)$ implica que el homomorfismo canónico
\[
  u:\Omega^1_k\otimes_kK\longrightarrow\Omega^1_A\otimes_AK
\]
es inyectivo. Ahora bien, se tiene para todo $\alpha$ un diagrama conmutativo
\[
\xymatrix@C=5em@R=3em{
  \Omega^1_k\otimes_kK \ar[r]^{u} \ar[d]_{v_\alpha} &
  \Omega^1_A\otimes_AK \ar[d] \\
  \Omega^1_{k/k_\alpha}\otimes_kK \ar[r]_{u_\alpha} &
  \Omega^1_{A/k_\alpha}\otimes_AK
}
\]

\oldpage[0\textsubscript{IV-1}]{190}

y la hipótesis de que $u_\alpha$ es inyectivo para
$\alpha\geqslant\alpha_0$ implica que $\operatorname{Ker}(u)$ está contenido
en la \emph{intersección} de los $\operatorname{Ker}(v_\alpha)$. Pero se ha
visto (21.8.3) que esta intersección es nula en virtud de la hipótesis
$\bigcap_\alpha k_\alpha(k^p)=k^p$, lo que termina la demostración.
\end{proof}

\begin{proposition}[22.2.9]
\label{0.22.2.9-es}
Bajo las hipótesis de (22.2.1), las condiciones siguientes son equivalentes:
\begin{enumerate}
  \item[a)] $A/\mathfrak m^2$ es una $k$-extensión $k$-trivial de $K$ por
  $\mathfrak m/\mathfrak m^2$.
  \item[b)] $\chi_{A/k}=0$.
  \item[c)] El homomorfismo canónico (20.5.11.2)
  \[
    \delta_{K/A/k}:\mathfrak m/\mathfrak m^2
    \longrightarrow\Omega^1_{A/k}\otimes_AK
  \]
  es inyectivo (en otras palabras, la sucesión
  \begin{equation}
  \tag{22.2.9.1}
  \label{0.22.2.9.1-es}
    0\longrightarrow\mathfrak m/\mathfrak m^2
    \longrightarrow\Omega^1_{A/k}\otimes_AK
    \longrightarrow\Omega^1_{K/k}\longrightarrow0
  \end{equation}
  es exacta).
\end{enumerate}
\end{proposition}

\begin{proof}
La equivalencia de $b)$ y $c)$ resulta de la sucesión exacta (20.6.22.1); la
equivalencia de $c)$ y $a)$ resulta de (20.5.12), pues la sucesión (22.2.9.1) se
escinde si es exacta, ya que se trata de espacios vectoriales.

Se observará que, si $K$ es \emph{separable sobre $k$}, se satisfacen las
condiciones equivalentes de (22.2.9), en virtud de la definición de
$\chi_{A/k}$ (22.2.1) y de (20.6.3).
\end{proof}

\begin{corollary}[22.2.10]
\label{0.22.2.10-es}
Bajo las hipótesis de (22.2.1), consideremos un subcuerpo $k_0$ de $k$. Las
condiciones siguientes son equivalentes:
\begin{enumerate}
  \item[a)] $A/\mathfrak m^2$ es una $k_0$-extensión $k_0$-trivial de $K$
  por $\mathfrak m/\mathfrak m^2$.
  \item[b)] El homomorfismo compuesto
  \[
    \Upsilon_{K/k_0}\longrightarrow\Upsilon_{K/k}
    \xrightarrow{\chi_{A/k}}\mathfrak m/\mathfrak m^2
  \]
  es nulo.
  \item[c)] El homomorfismo característico
  $\chi_{A/k}:\Upsilon_{K/k}\to\mathfrak m/\mathfrak m^2$ se factoriza como
  \[
    \Upsilon_{K/k}\longrightarrow\Upsilon_{K/k/k_0}
    \xrightarrow{\chi'}\mathfrak m/\mathfrak m^2.
  \]
\end{enumerate}
Además, cuando se cumplen estas condiciones, el homomorfismo $\chi'$ queda
determinado de manera única y coincide con el homomorfismo característico de
la $k_0$-extensión $k_0$-trivial $A/\mathfrak m^2$ de $K$ por
$\mathfrak m/\mathfrak m^2$.
\end{corollary}

\begin{proof}
Como el compuesto de $b)$ no es otro que $\chi_{A/k_0}$, la equivalencia de
$a)$ y $b)$ resulta de (22.2.9); por otra parte, $b)$ y $c)$ son equivalentes
en virtud de
la sucesión exacta (20.6.18.1)
\[
  \Upsilon_{K/k_0}\longrightarrow\Upsilon_{K/k}
  \longrightarrow\Upsilon_{K/k/k_0}\longrightarrow0,
\]
y, como el homomorfismo canónico
$\Upsilon_{K/k}\to\Upsilon_{K/k/k_0}$ es sobreyectivo, esto implica la
unicidad de $\chi'$; el hecho de que $\chi'$ sea el homomorfismo característico
de la $k_0$-extensión $A/\mathfrak m^2$ resulta de (20.6.22.2).
\end{proof}

\begin{corollary}[22.2.11]
\label{0.22.2.11-es}
Bajo las hipótesis de (22.2.1), sea $p$ el exponente característico de $k$ y
sea $(k_\alpha)$ una familia filtrante decreciente de subcuerpos de $k$ tal
que $\bigcap_\alpha k_\alpha(k^p)=k^p$. Si $\Upsilon_{K/k}$ es de rango finito
sobre $K$, existe un índice $\alpha$ tal que $A/\mathfrak m^2$ sea una
$k_\alpha$-extensión $k_\alpha$-trivial de $K$ por
$\mathfrak m/\mathfrak m^2$.
\end{corollary}

\begin{proof}
En efecto, se sabe (21.8.3) que existe entonces un índice $\alpha$ tal que el
homomorfismo canónico $\Upsilon_{K/k}\to\Upsilon_{K/k/k_\alpha}$ sea
biyectivo, luego se satisface la condición $b)$ de (22.2.10) (donde $k_\alpha$
sustituye a $k_0$).
\end{proof}

\oldpage[0\textsubscript{IV-1}]{191}

\subsection{Aplicación a las relaciones entre ciertos anillos locales y sus completados}
\label{subsection:0.22.3-es}

\begin{lemma}[22.3.1]
\label{0.22.3.1-es}
Sean $A_0$ un anillo semilocal noetheriano y $A$ una $A_0$-álgebra finita
(que es, por tanto, un anillo semilocal noetheriano, cf. Bourbaki,
\emph{Alg. comm.}, cap.~IV, §~2, n\textsuperscript{o}~5, cor.~3 de la
prop.~9). Si $\widehat A_0$ y $\widehat A$ son los completados de $A_0$ y $A$
para sus topologías preádicas respectivas, entonces, cuando se provee a
$A_0$, $A$ y $\widehat A$ de las topologías discretas, $\widehat A$ es una
$A$-álgebra formalmente lisa relativamente a $A_0$ (19.9.1).
\end{lemma}

\begin{proof}
En efecto, se sabe que
$\widehat A=A\otimes_{A_0}\widehat A_0$ (Bourbaki,
\emph{Alg. comm.}, cap.~IV, §~2, n\textsuperscript{o}~5, cor.~3 de la
prop.~9 y cap.~III, §~3, n\textsuperscript{o}~4, th.~1), luego la proposición
resulta de (19.9.3).
\end{proof}

\begin{proposition}[22.3.2]
\label{0.22.3.2-es}
Sean $A$ un anillo local noetheriano regular, $K$ su cuerpo de fracciones y
$p$ la característica de $K$. Se supone que se cumple una de las hipótesis
siguientes:
\begin{enumerate}
  \item[(i)] $p=0$.
  \item[(ii)] $p>0$ y existe una familia filtrante decreciente $(A_\alpha)$
  de subanillos noetherianos de $A$, tales que $A$ sea una
  $A_\alpha$-álgebra finita para todo $\alpha$, que
  $A^{p^{h(\alpha)}}\subset A_\alpha$ para un entero $h(\alpha)$ conveniente,
  y que, si $K_\alpha$ denota el cuerpo de fracciones de $A_\alpha$, se tenga
  \[
    \bigcap_\alpha K_\alpha(K^p)=K^p.
  \]
\end{enumerate}
Sean entonces $B$ una $A$-álgebra finita íntegra que contiene a $A$,
$\mathfrak n$ un ideal primo de $B$, $C$ el anillo local $B_\mathfrak n$ y
$\mathfrak q$ un ideal primo del completado $\widehat C$ tal que
$\mathfrak q\cap C=0$, de modo que el anillo local
$\widehat C_\mathfrak q$ es un álgebra sobre el cuerpo de fracciones $L$ de
$B$. Entonces $\widehat C_\mathfrak q$ es una $L$-álgebra formalmente lisa
para su topología $\mathfrak q$-ádica y, por consiguiente, un anillo
geométricamente regular sobre $L$ (19.6.5).
\end{proposition}

\begin{proof}
Distinguamos dos casos.

I) Supongamos que $\mathfrak n$ contiene el ideal maximal $\mathfrak m$ de
$A$, y por consiguiente $\mathfrak n\cap A=\mathfrak m$. Entonces
$\mathfrak n$ es maximal (Bourbaki, \emph{Alg. comm.}, cap.~V, §~2,
n\textsuperscript{o}~1, prop.~1) y, si $\mathfrak r$ es el radical del anillo
semilocal $B$, $\widehat C$ es el completado separado de $B$ para la topología
$\mathfrak n$-preádica, y uno de los componentes del anillo semilocal
$\widehat B$, completado de $B$ para la topología $\mathfrak r$-preádica
(Bourbaki, \emph{Alg. comm.}, cap.~III, §~3, n\textsuperscript{o}~4,
prop.~8); se puede escribir, por tanto,
$\widehat C_\mathfrak q=\widehat B_{\mathfrak q'}$, donde $\mathfrak q'$ es
el ideal primo de $\widehat B$ imagen recíproca de $\mathfrak q$. Observemos
ahora que $B$ es un subanillo de $\widehat B$ y la hipótesis
$\mathfrak q\cap C=0$ implica $\mathfrak q'\cap B=0$, luego
$\widehat B_{\mathfrak q'}$ es también una localización del anillo
$\widehat B\otimes_B L$ en uno de sus ideales primos $\mathfrak q''$, y
$\mathfrak q'$ es su imagen recíproca en $\widehat B$. Además, se tiene
$\widehat B=\widehat A\otimes_A B$ (Bourbaki, \emph{Alg. comm.}, cap.~IV,
§~2, n\textsuperscript{o}~5, cor. de la prop.~9 y cap.~III, §~3,
n\textsuperscript{o}~4, th.~1), luego
$\widehat B\otimes_B L=\widehat A\otimes_A L$; finalmente, si $\mathfrak p$
es el ideal primo de $\widehat A$, imagen recíproca de $\mathfrak q''$, el
hecho de que $A\to B$ sea inyectivo y la conmutatividad del diagrama
\[
  \xymatrix{
    B \ar[r] & \widehat B\\
    A \ar[u] \ar[r] & \widehat A \ar[u]
  }
\]
implican que $\mathfrak p\cap A=0$, luego
$(\widehat A\otimes_A L)_{\mathfrak q''}$ es también una localización del
anillo
\[
  \widehat A_\mathfrak p\otimes_A L
  =(\widehat A_\mathfrak p\otimes_A K)\otimes_K L
\]
en uno de sus ideales primos.

\oldpage[0\textsubscript{IV-1}]{192}

Para mostrar que este anillo local es una $L$-álgebra formalmente lisa, basta
probar que $\widehat A_\mathfrak p\otimes_A K$ es una $K$-álgebra formalmente
lisa (19.3.5, (iii) y (iv)); por otra parte,
$\widehat A_\mathfrak p\otimes_A K=\widehat A_\mathfrak p$, puesto que
$\mathfrak p\cap A=0$. Vamos a aplicar los criterios (22.2.2) y (22.2.8).
En primer lugar, como $A$ es regular, también lo es $\widehat A$ (17.1.5),
luego $\widehat A_\mathfrak p$ (17.3.2), y se cumple la condición (i) de
(22.2.2) (17.1.1). Basta por consiguiente (cuando $p>0$) verificar la
condición b) de (22.2.8). Ahora bien, (22.3.1) muestra que, para todo
$\alpha$, $\widehat A$ es una $A$-álgebra formalmente lisa relativamente a
$A_\alpha$ (para las topologías discretas); por consiguiente,
\[
  \widehat A\otimes_{A_\alpha}K_\alpha
  =\widehat A\otimes_A(A\otimes_{A_\alpha}K_\alpha)
\]
es una $(A\otimes_{A_\alpha}K_\alpha)$-álgebra formalmente lisa relativamente
a $K_\alpha$, para las topologías discretas (19.9.2, (iii)). Pero, como $A$
es una $A_\alpha$-álgebra finita, se tiene
$A\otimes_{A_\alpha}K_\alpha=K$; como $\widehat A_\mathfrak p$ es un anillo
de fracciones de $\widehat A\otimes_A K$, se concluye (19.9.2, (iv)) que
$\widehat A_\mathfrak p$ es una $K$-álgebra formalmente lisa relativamente a
$K_\alpha$, para la topología discreta, y \emph{a fortiori} para su topología
$\mathfrak p$-ádica ((19.3.8) y (19.9.5)). Pero esta es precisamente la
condición b) de (22.2.8), luego la proposición queda demostrada en este caso
para $p>0$. Cuando $p=0$, el cuerpo residual de
$\widehat A_\mathfrak p$ es separable sobre $K$, y como
$\widehat A_\mathfrak p$ es un anillo regular, es una $K$-álgebra formalmente
lisa en virtud de (19.6.4).

II) Caso general. Sea $\mathfrak s$ el ideal primo $\mathfrak n\cap A$ de
$A$, y pongamos $S=A-\mathfrak s$, de modo que $S^{-1}A=A_\mathfrak s$ y se
tiene $C=(S^{-1}B)_{S^{-1}\mathfrak n}$, donde ahora
$S^{-1}\mathfrak n$ contiene el ideal maximal
$\mathfrak sA_\mathfrak s$ de $A_\mathfrak s$. Como $A_\mathfrak s$ es
regular (17.3.2) y $S^{-1}B$ una $A_\mathfrak s$-álgebra finita íntegra que
contiene a $A_\mathfrak s$, todo se reduce a verificar la condición (ii) para
$A_\mathfrak s$ cuando $p>0$: ahora bien, pongamos
$\mathfrak s_\alpha=\mathfrak s\cap A_\alpha$ y consideremos el anillo local
$(A_\alpha)_{\mathfrak s_\alpha}$. Para todo $t\notin\mathfrak s$, por
hipótesis, poniendo $q=p^{h(\alpha)}$, se tiene $t^q\in A_\alpha$ y
$t^q\notin\mathfrak s$, luego $t^q\notin\mathfrak s_\alpha$; si se pone
$S_\alpha=A_\alpha-\mathfrak s_\alpha$, se ve por tanto que
$A_\mathfrak s=S_\alpha^{-1}A$, y la hipótesis implica que
$A_\mathfrak s$ es una $(A_\alpha)_{\mathfrak s_\alpha}$-álgebra finita;
además, $K_\alpha$ es también el cuerpo de fracciones de
$(A_\alpha)_{\mathfrak s_\alpha}$, lo que termina la demostración.
\end{proof}

\begin{theorem}[22.3.3]
\label{0.22.3.3-es}
Sean $A$ un anillo local noetheriano completo, $\mathfrak p$ un ideal primo
de $A$, $B$ el anillo localizado $A_\mathfrak p$, $\mathfrak q$ un ideal
primo del completado $\widehat B$ y $\mathfrak r=\mathfrak q\cap B$.
Entonces el anillo localizado $\widehat B_\mathfrak q$ de $\widehat B$ es una
$B_\mathfrak r$-álgebra formalmente lisa para las topologías preádicas.
\end{theorem}

\begin{proof}
El ideal primo $\mathfrak r$ de $B$ es de la forma
$\mathfrak n_\mathfrak p$, donde $\mathfrak n\subset\mathfrak p$ es un ideal
primo de $A$, y se tiene $B_\mathfrak r=A_\mathfrak n$; sea $k$ el cuerpo
residual de $B_\mathfrak r$, que es también el cuerpo de fracciones del anillo
local íntegro y completo $A'=A/\mathfrak n$. Como $\widehat B$ es un
$B$-módulo plano $(\mathbf 0_\mathrm I,~7.3.5)$,
$\widehat B_\mathfrak q$ es un $B_\mathfrak r$-módulo plano
$(\mathbf 0_\mathrm I,~6.3.2)$, y para probar que
$\widehat B_\mathfrak q$ es una $B_\mathfrak r$-álgebra formalmente lisa
basta probar que
\[
  \widehat B_\mathfrak q\otimes_{B_\mathfrak r}k
  =\widehat B_\mathfrak q/\mathfrak r\widehat B_\mathfrak q
\]
es una $k$-álgebra formalmente lisa (19.7.1). Ahora bien,
$\widehat B_\mathfrak q/\mathfrak r\widehat B_\mathfrak q
=(\widehat B/\mathfrak r\widehat B)_{\mathfrak q'}$, donde $\mathfrak q'$ es
la imagen canónica de $\mathfrak q$ en
$\widehat B/\mathfrak r\widehat B$; por definición se tiene
$\mathfrak q'\cap(B/\mathfrak r)=0$ y
$B/\mathfrak r=A_{\mathfrak p'}$, donde
$\mathfrak p'=\mathfrak p/\mathfrak n$; además,
$\widehat B/\mathfrak r\widehat B$ es el completado de $B/\mathfrak r$.

Quedamos así reducidos a demostrar el corolario siguiente:
\end{proof}

\begin{corollary}[22.3.4]
\label{0.22.3.4-es}
Sean $A$ un anillo local noetheriano completo e íntegro, $K$ su cuerpo de
fracciones, $\mathfrak p$ un ideal primo de $A$ y $B$ el anillo localizado
$A_\mathfrak p$. Entonces, para todo ideal primo $\mathfrak q$

\oldpage[0\textsubscript{IV-1}]{193}

del completado $\widehat B$, tal que $\mathfrak q\cap B=0$, el anillo local
$\widehat B_\mathfrak q$ es una $K$-álgebra formalmente lisa para su topología
$\mathfrak q$-ádica y, por tanto, un anillo geométricamente regular sobre
$K$.
\end{corollary}

\begin{proof}
En efecto, resulta de (19.8.9) que existe un subanillo $A_0$ de $A$ que es un
anillo de series formales sobre un cuerpo de característica $p>0$ o sobre un
anillo de Cohen (se recuerda que los anillos de Cohen tienen un cuerpo de
fracciones de característica $0$), y tal que $A$ es una $A_0$-álgebra finita.
Ahora bien, se sabe que el anillo $A_0$ es regular (17.3.8) y cumple una de las
hipótesis (i), (ii) de (22.3.2), en virtud de (21.8.8); basta entonces aplicar
(22.3.2), sustituyendo $B$ por $A$ y $A$ por $A_0$.
\end{proof}
\subsection{Resultados preliminares sobre las extensiones finitas de anillos locales cuyo ideal maximal es de cuadrado nulo}
\label{subsection:0.22.4-es}
\label{0.22.4-es}

\begin{proposition}[22.4.1]
\label{0.22.4.1-es}
Sean $A$ un anillo local cuyo ideal maximal $\mathfrak m$ es de cuadrado
nulo, $K=A/\mathfrak m$ su cuerpo residual; denotemos por $V$ el ideal
$\mathfrak m$ considerado como $K$-espacio vectorial. Sean $T_i$
$(1\leqslant i\leqslant r)$ indeterminadas y, para cada $i$, sea
$F_i$ un polinomio mónico de $A[T_i]$ de grado $d_i$; denotemos por $A'$ el
cociente del anillo de polinomios $A[T_1,\ldots,T_r]$ por el ideal
$\mathfrak b$ engendrado por los $r$ polinomios $F_i$. Supongamos que $A'$ es
un anillo local, y sean $\mathfrak m'$ su ideal maximal y
$K'=A'/\mathfrak m'$ su cuerpo residual. Entonces:
\begin{enumerate}
  \item[(i)] Si se pone $d=\prod_{i=1}^r d_i$, $A'$ es un $A$-módulo
  libre de rango $d$, y se tiene $[K':K]\leqslant d$.
  \item[(ii)] Para que $[K':K]=d$, es necesario y suficiente que
  $A'\otimes_AK=A'/\mathfrak mA'$ sea un cuerpo (necesariamente isomorfo a
  $K'$), es decir, que $\mathfrak m'=\mathfrak mA'$. En este caso se tiene
  $\mathfrak m'^2=0$, y si $V'$ denota el ideal $\mathfrak m'$ considerado
  como $K'$-espacio vectorial, el homomorfismo canónico
  $V\otimes_KK'\to V'$ es biyectivo (y, por consiguiente,
  $\operatorname{rg}_{K'}V'=\operatorname{rg}_K V$).
\end{enumerate}
\end{proposition}

\begin{proof}
Sin suponer que $A'$ sea local, es evidente, por división euclídea e
inducción sobre $r$, que si $t_i$ es la clase de $T_i$ mod.~$\mathfrak b$,
los monomios
$M_\nu=t_1^{\nu_1}t_2^{\nu_2}\cdots t_r^{\nu_r}$, donde
$0\leqslant\nu_i\leqslant d_i-1$, forman una base del $A$-módulo $A'$. Si
$A'$ se supone local, $K'$ es un cociente de
$A'/\mathfrak mA'=A'\otimes_AK$, por tanto
$[K':K]\leqslant[A'\otimes_AK:K]=d$, y no puede haber igualdad salvo si
$\mathfrak m'=\mathfrak mA'$, lo que demuestra (i) y la primera afirmación
de (ii). Por otra parte, está claro que cuando
$\mathfrak m'=\mathfrak mA'$, se tiene $\mathfrak m'^2=0$, y
$V\otimes_KK'=\mathfrak m\otimes_AA'=\mathfrak mA'=V'$ puesto que $A'$ es
un $A$-módulo libre.
\end{proof}

\begin{remark}[22.4.2]
\label{0.22.4.2-es}
Cuando $A$ es artiniano (es decir, $\operatorname{rg}_K(V)$ es finito),
la hipótesis de que $A'$ sea local implica siempre que
$\operatorname{rg}_{K'}(\mathfrak m'/\mathfrak m'^2)
\geqslant\operatorname{rg}_K(V)$, como veremos más adelante (22.5.2).
Obsérvese que estos dos rangos pueden ser iguales sin que se tenga
$\mathfrak m'=\mathfrak mA'$.
\end{remark}

\begin{proposition}[22.4.3]
\label{0.22.4.3-es}
Siendo las hipótesis sobre $A$ y las notaciones las de (22.4.1), supongamos
que los $F_i$ son de la forma
\begin{equation}
\label{0.22.4.3.1-es}
  F_i(T_i)=T_i^{d_i}+\sum_{1\leqslant k\leqslant d_i}
  c_{ik}T_i^{d_i-k},
  \tag{22.4.3.1}
\end{equation}
donde $d_i\geqslant2$ para todo $i$, y $c_{ik}\in\mathfrak m$
(«polinomios de Eisenstein»). Entonces:
\begin{enumerate}
  \item[(i)] $A'$ es un anillo local, y su cuerpo residual
  $K'=A'/\mathfrak m'$ es isomorfo a $K$.

  \oldpage[0\textsubscript{IV-1}]{194}

  \item[(ii)] El núcleo del homomorfismo canónico
  $V\to V'=\mathfrak m'/\mathfrak m'^2$ es el ideal $\mathfrak n$ de $A$
  engendrado por los $\xi_i=c_{i,d_i}$ («términos constantes» de los $F_i$), y
  el conúcleo de este homomorfismo tiene por base sobre $K'$ las imágenes de los $T_i$;
  en particular, se tiene
  \begin{equation}
  \label{0.22.4.3.2-es}
    \operatorname{rg}_{K'}(V')=\operatorname{rg}_K(V)+r-r',
    \tag{22.4.3.2}
  \end{equation}
  donde $r'=\operatorname{rg}_K(\mathfrak n)$. Cuando
  $\operatorname{rg}_K(V)$ es finito (es decir, cuando $A$ es
  artiniano), para que
  $\operatorname{rg}_{K'}(V')=\operatorname{rg}_K(V)$, es necesario y suficiente
  que los $\xi_i$ sean linealmente independientes sobre $K$.
\end{enumerate}
\end{proposition}

\begin{proof}
(i) Está claro que $A'\otimes_AK=A'/\mathfrak mA'$ es isomorfo a
$K[T_1,\ldots,T_r]/\mathfrak b'$, donde $\mathfrak b'$ es el ideal engendrado
por los $r$ polinomios $T_i^{d_i}$; es, pues, un anillo local de cuerpo
residual $K$, cuyo ideal maximal está engendrado por las clases de los $T_i$
mod.~$\mathfrak b'$, de donde (i).

(ii) Sea $t_i$ $(1\leqslant i\leqslant r)$ la clase de $T_i$
mod.~$\mathfrak b$; se deduce de (i) que el ideal maximal $\mathfrak m'$ de
$A'$ viene dado por
\[
  \mathfrak m'=\mathfrak mA'+\sum_{i=1}^r t_iA',
\]
de donde, teniendo en cuenta que $\mathfrak m^2=0$,
\[
  \mathfrak m'^2=\sum_{i,j}t_it_jA'+\sum_i t_i\mathfrak mA'.
\]
Se concluye que $A'/\mathfrak m'^2$ es isomorfo al anillo
$A[T_1,\ldots,T_r]/\mathfrak c$, donde $\mathfrak c$ es el ideal engendrado por
los elementos
\[
  F_i\ (1\leqslant i\leqslant r),\qquad
  T_iT_j\ (1\leqslant i\leqslant r,\ 1\leqslant j\leqslant r),\qquad
  xT_i\ (1\leqslant i\leqslant r,\ x\in\mathfrak m).
\]
Como $d_i\geqslant2$ y los $T_i^2$ pertenecen a
$\mathfrak c$, las hipótesis hechas implican que $\mathfrak c$ está también
engendrado por los elementos
\[
  \xi_i\ (1\leqslant i\leqslant r),\qquad
  T_iT_j\ (1\leqslant i\leqslant r,\ 1\leqslant j\leqslant r),\qquad
  T_ix\ (1\leqslant i\leqslant r,\ x\in\mathfrak m).
\]
Sean
\[
  \mathfrak c_1=\sum_{i,j}T_iT_jA[T_1,\ldots,T_r],\qquad
  \mathfrak c_2=\mathfrak c_1+\sum_i\mathfrak mT_iA[T_1,\ldots,T_r],
\]
de modo que
$\mathfrak c=\mathfrak c_2+\sum_i\xi_iA[T_1,\ldots,T_r]$. Está claro que
$A[T_1,\ldots,T_r]/\mathfrak c_1$ es el $A$-módulo suma directa de $A$ y
de los $At_i^{(1)}$, donde $t_i^{(1)}$ es la clase de $T_i$
mod.~$\mathfrak c_1$. Se deduce que
$A[T_1,\ldots,T_r]/\mathfrak c_2$ es suma directa de $A$ y de los
$Kt_i^{(2)}$, donde $t_i^{(2)}$ es la clase de $T_i$
mod.~$\mathfrak c_2$. Finalmente, $A[T_1,\ldots,T_r]/\mathfrak c$ es suma
directa de $A/\mathfrak n$ y de los $Kt_i'$, donde $t_i'$ es la clase de $T_i$
mod.~$\mathfrak c$, y $\mathfrak n$ es el ideal de $A$ engendrado por los
$\xi_i$; en la identificación de $A'/\mathfrak m'^2$ con
$A[T_1,\ldots,T_r]/\mathfrak c$, la imagen de $V$ se identifica con
$\mathfrak m/\mathfrak n$, y lo anterior demuestra, por tanto, (ii).
\end{proof}

\begin{proposition}[22.4.4]
\label{0.22.4.4-es}
Siendo las hipótesis sobre $A$ y las notaciones las de (22.4.1), supongamos
que $K$ es de característica $p>0$ (por tanto $p\cdot1\in\mathfrak m$ en
$A$) y que los $F_i$ son de la forma
\begin{equation}
\label{0.22.4.4.1-es}
  F_i(T_i)=T_i^p+p\sum_{k=1}^{p-1}c_{ik}T_i^{p-k}-f_i
  \tag{22.4.4.1}
\end{equation}
con $c_{ik}\in A$ y $f_i\in A$ para todo $i$ y todo $k$ tal que
$1\leqslant k\leqslant p-1$; además, si $a_i$ es la clase de $f_i$
mod.~$\mathfrak m$, se supone que la familia $(a_i)_{1\leqslant i\leqslant r}$
es $p$-libre sobre $K^p$ («absolutamente $p$-libre»). Entonces:
\begin{enumerate}
  \item[(i)] $A'$ es un anillo local, de ideal maximal $\mathfrak mA'$,
  cuyo cuerpo residual $K'$ es isomorfo a la extensión de $K$ obtenida por
  adjunción de los $a_i^{1/p}$ $(1\leqslant i\leqslant r)$.

  \oldpage[0\textsubscript{IV-1}]{195}

  \item[(ii)] Los elementos $d_Af_i\otimes1_{K'}$ forman una base del núcleo
  $N=\Upsilon^{K'}_{A'/A}$ del homomorfismo canónico
  \[
    \Omega^1_A\otimes_AK'\longrightarrow
    \Omega^1_{A'}\otimes_{A'}K'.
  \]
\end{enumerate}
\end{proposition}

\begin{proof}
(i) Esta vez $A'\otimes_AK=A'/\mathfrak mA'$ es isomorfo a
$K[T_1,\ldots,T_r]/\mathfrak b'$, donde $\mathfrak b'$ es el ideal engendrado
por los polinomios $T_i^p-a_i$; la hipótesis hecha sobre los $a_i$ implica que
este anillo cociente es un cuerpo (21.1.8), de donde las afirmaciones de (i).

Antes de demostrar (ii), estableceremos el

\begin{lemma}[22.4.4.2]
\label{0.22.4.4.2-es}
Sean $\Lambda$ un anillo, $A$ una $\Lambda$-álgebra y $\mathfrak J$ un
ideal de $A$ tal que el anillo $C=A/\mathfrak J$ sea de característica
$p>0$ (lo que equivale a decir que $p\cdot1\in\mathfrak J$ en $A$). Sea
$\pi$ la imagen canónica de $p\cdot1$ en
$\mathfrak J/\mathfrak J^2$. Si $C$ es una $(\Lambda/p\Lambda)$-álgebra
formalmente lisa para las topologías discretas, se tiene una sucesión exacta
canónica
\begin{equation}
\label{0.22.4.4.3-es}
  0\longrightarrow(\mathfrak J/\mathfrak J^2)/C\cdot\pi
  \longrightarrow\Omega^1_{A/\Lambda}\otimes_AC
  \longrightarrow\Omega^1_{C/\Lambda}\longrightarrow0.
  \tag{22.4.4.3}
\end{equation}
\end{lemma}

\begin{proof}[Demostración del lema]
Pongamos, en efecto, $A'=A/pA$, $\mathfrak J'=\mathfrak J/pA$, de modo que
$C=A'/\mathfrak J'$. Como
$A'=A\otimes_\Lambda(\Lambda/p\Lambda)$, se tiene
$\Omega^1_{A/\Lambda}\otimes_AA'=\Omega^1_{A'/\Lambda'}$ salvo
isomorfismo canónico, donde $\Lambda'$ denota el anillo
$\Lambda/p\Lambda$ (20.5.5). Aplicando entonces la sucesión exacta (20.5.14.1) a
la $\Lambda'$-álgebra formalmente lisa $C=A'/\mathfrak J'$, se obtiene la
sucesión exacta
\[
  0\longrightarrow\mathfrak J'/\mathfrak J'^2
  \longrightarrow\Omega^1_{A'/\Lambda'}\otimes_{A'}C
  \longrightarrow\Omega^1_{C/\Lambda'}\longrightarrow0
\]
y como el homomorfismo $\Lambda\to\Lambda'$ es sobreyectivo, se tiene
$\Omega^1_{C/\Lambda'}=\Omega^1_{C/\Lambda}$ salvo isomorfismo canónico
(20.5.15); de donde la sucesión exacta (22.4.4.3) en virtud de lo anterior.
\end{proof}

(ii) Como se tiene $\mathfrak m'^2=0$, denotaremos también por $V'$ el ideal
$\mathfrak m'$ considerado como $K'$-espacio vectorial, y pondremos
\begin{equation}
\label{0.22.4.4.4-es}
  V_0=V/K\cdot p=\mathfrak m/(\mathfrak m^2+pA),\qquad
  V_0'=V'/K'\cdot p=\mathfrak m'/(\mathfrak m'^2+pA').
  \tag{22.4.4.4}
\end{equation}
Aplicando la sucesión exacta (22.4.4.3) al caso en que $\Lambda=\mathbf Z$, a la
$\mathbf Z$-álgebra $A$ (resp. $A'$) y a su ideal $\mathfrak m$
(resp. $\mathfrak m'$), se obtienen, puesto que $K$ (resp. $K'$) es separable sobre
el cuerpo primo $\mathbf Z/p\mathbf Z$, y por tanto una
$(\mathbf Z/p\mathbf Z)$-álgebra formalmente lisa (19.6.1), las sucesiones
exactas
\begin{equation}
\label{0.22.4.4.5-es}
  0\longrightarrow V_0\longrightarrow\Omega^1_A\otimes_AK
  \longrightarrow\Omega^1_K\longrightarrow0,\qquad
  0\longrightarrow V_0'\longrightarrow\Omega^1_{A'}\otimes_{A'}K'
  \longrightarrow\Omega^1_{K'}\longrightarrow0.
  \tag{22.4.4.5}
\end{equation}
Consideremos entonces el diagrama (22.4.4.6), donde las dos columnas centrales son
la segunda sucesión (22.4.4.5) y la primera tensorizada por $K'$; como la
primera sucesión exacta (22.4.4.5) está formada por espacios vectoriales sobre $K$,
las dos columnas centrales de (22.4.4.6) son exactas; además, el
diagrama formado por estas columnas y las flechas horizontales que las
conectan es conmutativo, en virtud de la demostración de (22.4.4.2) y de
la conmutatividad del diagrama (20.5.11.3). La última fila de (22.4.4.6)
es la sucesión exacta obtenida a partir de (20.6.1.1) sustituyendo $A$, $B$,
$C$ por $\mathbf Z$, $K$ y $K'$; la fila central se obtiene por
tensorización con $K'$ de la sucesión exacta de $K$-módulos deducida de
(20.6.1.1) sustituyendo $A$, $B$, $C$ por $\mathbf Z$, $A$, $K$
respectivamente. El diagrama formado por estas dos filas y las flechas
verticales que las conectan es conmutativo en virtud de la conmutatividad de
(20.6.4.3). Observemos, por otra parte, que se está en las condiciones de aplicación
de (22.4.1, (ii)), por lo que la fila superior del diagrama (22.4.4.6) es
exacta, y las columnas extremas del diagrama

\oldpage[0\textsubscript{IV-1}]{196}

\begin{equation}
\label{0.22.4.4.6-es}
  \xymatrix{
    && 0 \ar[d] & 0 \ar[d] && \\
    & 0 \ar[r] & V_0\otimes_KK' \ar[r] \ar[d] & V_0' \ar[r] \ar[d] & 0 & \\
    0 \ar[r] & N \ar[r] \ar[d] & \Omega^1_A\otimes_AK' \ar[r] \ar[d] &
      \Omega^1_{A'}\otimes_{A'}K' \ar[r] \ar[d] &
      \Omega^1_{A'/A}\otimes_{A'}K' \ar[r] \ar[d] & 0 \\
    0 \ar[r] & \Upsilon_{K'/K} \ar[r] \ar[d] &
      \Omega^1_K\otimes_KK' \ar[r] \ar[d] & \Omega^1_{K'} \ar[r] \ar[d] &
      \Omega^1_{K'/K} \ar[r] \ar[d] & 0 \\
    & 0 & 0 & 0 & 0 &
  }
  \tag{22.4.4.6}
\end{equation}
están formadas, respectivamente, por los núcleos y los conúcleos de las flechas
horizontales centrales, lo que termina de probar que el diagrama es
conmutativo. Además:

\begin{lemma}[22.4.4.7]
\label{0.22.4.4.7-es}
Los homomorfismos
\[
  N\longrightarrow\Upsilon_{K'/K}\qquad\text{y}\qquad
  \Omega^1_{A'/A}\otimes_{A'}K'\longrightarrow\Omega^1_{K'/K}
\]
del diagrama (22.4.4.6) son biyectivos.
\end{lemma}

\begin{proof}
En efecto, el diagrama de la serpiente (Bourbaki, \emph{Alg. conmut.}, cap.~I,
§~1, n\textsuperscript{o}~4, prop.~2), aplicado a las dos columnas centrales de
(22.4.4.6), da una sucesión exacta
\[
  0\longrightarrow N\longrightarrow\Upsilon_{K'/K}\longrightarrow0
  \longrightarrow\Omega^1_{A'/A}\otimes_{A'}K'
  \longrightarrow\Omega^1_{K'/K}\longrightarrow0.
\]
Ahora bien, si $t_i$ es la imagen canónica de $T_i$ en $A'$, la relación
$F_i(t_i)=0$, junto con el hecho de que en $A'$ se tiene
$p\cdot1\in\mathfrak mA'\subset\mathfrak m'$, muestra inmediatamente que
$d_{A'}f_i\in\mathfrak m'\Omega^1_{A'}$, por tanto
$d_{A'}f_i\otimes1_{K'}=0$; esto demuestra que los
$d_Af_i\otimes1_{K'}$ pertenecen a $N$. Además, sus imágenes en
$\Upsilon_{K'/K}$ son los $d_Ka_i\otimes1_{K'}$, que forman una base de
$\Upsilon_{K'/K}$ sobre $K'$ (21.4.7); teniendo en cuenta (22.4.4.7), esto termina
de demostrar (22.4.4).
\end{proof}
\end{proof}
\begin{env}[22.4.5]
\label{0.22.4.5-es}
Consideremos ahora dos anillos $A$, $B$ \emph{de característica} $p>0$, y
dos homomorfismos
\[
  A\xrightarrow{i}B\xrightarrow{j}A
\]

\oldpage[0\textsubscript{IV-1}]{197}

que verifican las condiciones de (21.3.1). Sea además $\mathfrak m$ un ideal
\emph{de cuadrado nulo} en $A$; pongamos
\[
  K=A/\mathfrak m,\qquad
  B_{(K)}=B\otimes_AK=B/Bi(\mathfrak m)
\]
y sean
\[
  \varphi:A\longrightarrow K,\qquad
  \psi:B\longrightarrow B_{(K)}
\]
los homomorfismos canónicos.

Como $\varphi$ es sobreyectivo, $\Omega^1_{B_{(K)}/A}$ se identifica
canónicamente con $\Omega^1_{B_{(K)}/K}$ (20.5.15). Observemos que, puesto que
$p\geqslant2$, para todo $x\in\mathfrak m$, se tiene $j(i(x))=x^p=0$, es decir,
$j(i(\mathfrak m))=0$; pasando a los cocientes, se deduce, pues, de $j$ un
homomorfismo
\[
  j':B_{(K)}\longrightarrow A
\]
tal que, si $i'=\psi\circ i:A\longrightarrow B_{(K)}$, se tiene
\[
  j'(i'(x))=x^p\qquad\text{y}\qquad i'(j'(y))=y^p
  \qquad\text{para }x\in A\text{ e }y\in B_{(K)}.
\]
Poniendo, según las notaciones de (21.3.2),
\[
  \Theta_{B_{(K)}/A}
  =\Omega^1_{B_{(K)}/A}\otimes_{B_{(K)}}A_{[j']}
  =\Omega^1_{B_{(K)}/K}\otimes_{B_{(K)}}A_{[j']},
\]
se tiene, pues, (21.3.2.2) un homomorfismo canónico
\begin{equation}
\label{0.22.4.5.1-es}
  \pi_{B_{(K)}/A}:\Theta_{B_{(K)}/A}\longrightarrow\Omega^1_A.
  \tag{22.4.5.1}
\end{equation}

Por otra parte, se deducen de $i$ y $j$, pasando a los cocientes (teniendo en
cuenta que $j(i(\mathfrak m))=0$), homomorfismos
\[
  K\xrightarrow{i_0}B_{(K)}\xrightarrow{j_0}K
\]
de modo que se tiene el diagrama conmutativo
\[
  \xymatrix{
    A \ar[r]^i \ar[d]_\varphi & B \ar[r]^j \ar[d]^\psi & A \ar[d]^\varphi \\
    K \ar[r]_{i_0} & B_{(K)} \ar[r]_{j_0} \ar[ur]_{j'} & K
  }
\]
Está claro, además, que se tiene
\[
  j_0(i_0(s))=s^p\qquad\text{y}\qquad i_0(j_0(t))=t^p
  \qquad\text{para }s\in K\text{ y }t\in B_{(K)}.
\]
Se define, pues, del mismo modo
\[
  \Theta_{B_{(K)}/K}
  =\Omega^1_{B_{(K)}/K}\otimes_{B_{(K)}}K_{[j_0]}
  =\Theta_{B_{(K)}/A}\otimes_AK,
\]
por lo que, tensorizando (22.4.5.1) por $K$, se obtiene un
$K$-\emph{homomorfismo canónico}
\begin{equation}
\label{0.22.4.5.2-es}
  \pi'_{B_{(K)}/K}=\pi_{B_{(K)}/A}\otimes1_K:
  \Theta_{B_{(K)}/K}\longrightarrow\Omega^1_A\otimes_AK.
  \tag{22.4.5.2}
\end{equation}

\oldpage[0\textsubscript{IV-1}]{198}

Se deduce inmediatamente de las definiciones que el diagrama
\[
  \xymatrix{
    \Theta_{B_{(K)}/K} \ar[r]^{\pi'_{B_{(K)}/K}} \ar[dr]_{\pi_{B_{(K)}/K}}
      & \Omega^1_A\otimes_AK \ar[d]^{\varphi_{K/A}} \\
    & \Omega^1_K
  }
\]
es conmutativo, siendo $\pi_{B_{(K)}/K}$ el homomorfismo canónico
correspondiente a $i_0$ y $j_0$ (21.3.2.2). Por \emph{restricción} a
$\Xi_{B_{(K)}/K}=\operatorname{Ker}(\pi_{B_{(K)}/K})$, el homomorfismo
$\pi'_{B_{(K)}/K}$ define, pues, un \emph{homomorfismo canónico}
\begin{equation}
\label{0.22.4.5.3-es}
  \Xi_{B_{(K)}/K}\longrightarrow
  \operatorname{Ker}(\Omega^1_A\otimes_AK\longrightarrow\Omega^1_K)=\Upsilon_{K/A}.
  \tag{22.4.5.3}
\end{equation}
\end{env}

\begin{env}[22.4.6]
\label{0.22.4.6-es}
Siendo las hipótesis y notaciones las de (22.4.5), supongamos además que
$\mathfrak m$ es un ideal \emph{maximal de cuadrado nulo}, de modo que $K$ es un
\emph{cuerpo}, y denotemos por $V$ el ideal $\mathfrak m$ considerado como
$K$-espacio vectorial. Como $K$ es un álgebra formalmente lisa sobre su cuerpo primo
$\mathbf F_p$ (21.5.2) y $A$ es un álgebra sobre $\mathbf F_p$ (de modo
que $\Omega^1_A=\Omega^1_{A/\mathbf F_p}$), se deduce de (20.5.14) que se tiene
una \emph{sucesión exacta}
\begin{equation}
\label{0.22.4.6.1-es}
  0\longrightarrow V\longrightarrow\Omega^1_A\otimes_AK
  \longrightarrow\Omega^1_K\longrightarrow0.
  \tag{22.4.6.1}
\end{equation}
Se tiene, pues, por (22.4.5.3) un homomorfismo denotado por
\begin{equation}
\label{0.22.4.6.2-es}
  \chi'_{B/A}:\Xi_{B_{(K)}/K}\longrightarrow V
  \tag{22.4.6.2}
\end{equation}
que llamaremos \emph{homomorfismo característico} de la $A$-álgebra
$B$ (relativo a $i$, $j$ y $\mathfrak m$). Se deduce de las definiciones que para
que $\chi'_{B/A}$ sea \emph{inyectivo}, es necesario y suficiente que
$\pi'_{B_{(K)}/K}$ sea inyectivo, pues el núcleo de este último está evidentemente
contenido en $\Xi_{B_{(K)}/K}$.
\end{env}

\begin{proposition}[22.4.7]
\label{0.22.4.7-es}
Sean $A$ un anillo local artiniano cuyo ideal maximal $\mathfrak m$ es de
cuadrado nulo, $K=A/\mathfrak m$ su cuerpo residual y $V$ el ideal $\mathfrak m$
considerado como $K$-espacio vectorial; se supone que $A$ es de
característica $p>0$; sean
\begin{equation}
\label{0.22.4.7.1-es}
  F_i(T_i)=T_i^p-f_i\qquad\text{con }f_i\in A\quad(1\leqslant i\leqslant r)
  \tag{22.4.7.1}
\end{equation}
y denotemos por $B$ el anillo cociente de $A[T_1,\ldots,T_r]$ por el ideal
engendrado por los $r$ polinomios $F_i$. Entonces:
\begin{enumerate}
  \item[\rm(i)] $B$ es un anillo local.
  \item[\rm(ii)] Si $\mathfrak m'$ es el ideal maximal de $B$,
  $K'=B/\mathfrak m'$ su cuerpo residual y $V'$ el $B$-módulo
  $\mathfrak m'/\mathfrak m'^2$ considerado como $K'$-espacio vectorial, entonces,
  para que $\operatorname{rg}_K(V)=\operatorname{rg}_{K'}(V')$, es necesario y suficiente que el homomorfismo
  característico de $K$-espacios vectoriales
  \begin{equation}
  \label{0.22.4.7.2-es}
    \chi'_{B/A}:\Xi_{B_{(K)}/K}\longrightarrow V
    \tag{22.4.7.2}
  \end{equation}
  (cf. (22.4.6.2)) sea inyectivo.
\end{enumerate}
\end{proposition}

\begin{proof}
Denotemos también por $a_i$ $(1\leqslant i\leqslant r)$ la clase de $f_i$
mod.~$\mathfrak m$; los $a_i$ ya no son necesariamente $p$-independientes
sobre $K^p$, pero siempre se puede suponer que la subfamilia

\oldpage[0\textsubscript{IV-1}]{199}

$(a_i)_{s+1\leqslant i\leqslant r}$ es $p$-libre sobre $K^p$ y forma una
$p$-base sobre $K^p$ del cuerpo $K^p(a_1,\ldots,a_r)$, de modo que se tiene
\[
  a_i\in K^p(a_{s+1},\ldots,a_r)\qquad\text{para }1\leqslant i\leqslant s.
\]
Denotemos por $A''$ el anillo cociente de $A[T_{s+1},\ldots,T_r]$ por el ideal
engendrado por los $F_i$ de índice $i\geqslant s+1$; entonces, por (22.4.4), $A''$
es un anillo local cuyo ideal maximal es $\mathfrak m''=\mathfrak mA''$
(por tanto de cuadrado nulo) y cuyo cuerpo residual es
$K''=K(a_{s+1}^{1/p},\ldots,a_r^{1/p})$. Se tiene, pues,
$a_i\in K''^p$ para $1\leqslant i\leqslant s$ y, por consiguiente, en $A''$,
el elemento $f_i$, para $1\leqslant i\leqslant s$, puede escribirse
\begin{equation}
\label{0.22.4.7.3-es}
  f_i=g_i^p+h_i
  \tag{22.4.7.3}
\end{equation}
donde $g_i\in A''$ y $h_i\in\mathfrak m''$ (como $\mathfrak m''^2=0$ y
$p\mathfrak m''=0$, es inmediato que $h_i$ está determinado de manera
\emph{única} por estas condiciones). Sustituyendo $T_i$ por $T_i+g_i$, se ve, pues,
que si se pone
\[
  G_i(T_i)=T_i^p+\sum_{k=1}^{p-1}\binom pk g_i^{p-k}T_i^k-h_i
  \qquad(1\leqslant i\leqslant s),
\]
$B$ es el anillo cociente de $A''[T_1,\ldots,T_s]$ por el ideal engendrado por los
$G_i$; ahora bien, todos los coeficientes de $G_i$, salvo el coeficiente dominante,
pertenecen a $\mathfrak m''$, por lo que se está en la situación de (22.4.3),
$B$ es un anillo local, y su cuerpo residual $K'$ es \emph{isomorfo} a
$K''$, lo que demuestra (i) (que, por lo demás, resulta directamente de que
$B^p\subset A$, y por tanto solo hay un ideal de $B$ por encima de
$\mathfrak m$).

Observemos ahora que para que (22.4.7.2) sea inyectivo, es necesario y suficiente
que lo sea
\[
  \pi'_{B_{(K)}/K}:\Omega^1_{B_{(K)}/K}\otimes_{B_{(K)}}K
  \longrightarrow\Omega^1_A\otimes_AK
\]
(por (22.4.6)). Ahora bien, $B_{(K)}=B/\mathfrak mB$ es el álgebra cociente de
$C=K[T_1,\ldots,T_r]$ por el ideal engendrado por los polinomios
$\overline F_i=T_i^p-a_i$, y como $d_{C/K}\overline F_i=0$, se deduce de
(20.5.13) que $\Omega^1_{B_{(K)}/K}$ es un $B_{(K)}$-módulo libre que tiene por
\emph{base} los $d_{B_{(K)}/K}t_i$, donde $t_i$ es la imagen canónica de $T_i$
$(1\leqslant i\leqslant r)$. Pero, por definición (22.4.5), la imagen canónica por
$j':B_{(K)}\longrightarrow A$ de una clase mod.~$\mathfrak mB$ de un elemento
$x\in B$ es el elemento $x^p\in A$; se deduce inmediatamente, puesto que $t_i^p=f_i$,
que la imagen de $d_{B_{(K)}/K}t_i\otimes1$ por $\pi'_{B_{(K)}/K}$ es
$d_Af_i\otimes1$ en $\Omega^1_A\otimes_AK$; la condición de que
$\pi'_{B_{(K)}/K}$ sea inyectivo equivale, pues, a que los
$d_Af_i\otimes1_K$ sean \emph{linealmente independientes sobre $K$}, o también
a que sus imágenes $d_Af_i\otimes1_{K'}$ sean \emph{linealmente independientes
sobre $K'$} en $\Omega^1_A\otimes_AK'$.

Ahora bien, si se aplica (22.4.4) a la $A$-álgebra $A''$, se ve que el núcleo del
homomorfismo canónico
$\Omega^1_A\otimes_AK'\longrightarrow\Omega^1_{A''}\otimes_{A''}K'$ tiene por
base los $d_Af_i\otimes1_{K'}$ para $s+1\leqslant i\leqslant r$. La condición
precedente equivale, pues, también a decir que los $d_{A''}f_i\otimes1_{K'}$ son
linealmente independientes en $\Omega^1_{A''}\otimes_{A''}K'$ para
$1\leqslant i\leqslant s$. Ahora bien, se tiene $d_{A''}f_i=d_{A''}h_i$ por (22.4.7.3);
por otra parte, si $V''$ es el $K'$-espacio vectorial $\mathfrak m''$, se tiene
$\operatorname{rg}_{K'}(V'')=\operatorname{rg}_K(V)$ (22.4.1, (ii)), y la condición
$\operatorname{rg}_{K'}(V')=\operatorname{rg}_K(V)$ equivale, pues, a
$\operatorname{rg}_{K'}(V'')=\operatorname{rg}_{K'}(V')$. Pero se deduce de (22.4.3) que esta última
relación equivale a que las clases $h_i$

\oldpage[0\textsubscript{IV-1}]{200}

en $V''$ sean \emph{linealmente independientes sobre $K'$}. Ahora bien, en
$\Omega^1_{A''}\otimes_{A''}K'$, los $d_{A''}h_i\otimes1_{K'}$ son las imágenes
canónicas (20.5.11.2) de los $h_i$ por la inyección
$V''\longrightarrow\Omega^1_{A''}\otimes_{A''}K'$ (cf. (22.4.6.1)), y esto
termina de demostrar (22.4.7).
\end{proof}

Además, se ha demostrado:

\begin{corollary}[22.4.7.4]
\label{0.22.4.7.4-es}
Para que $\operatorname{rg}_K(V)=\operatorname{rg}_{K'}(V')$, es necesario y suficiente que los elementos
$d_Af_i\otimes1_K$ $(1\leqslant i\leqslant r)$ sean linealmente
independientes en el $K$-espacio vectorial $\Omega^1_A\otimes_AK$.
\end{corollary}

\begin{remark}[22.4.8]
\label{0.22.4.8-es}
Se pueden definir los homomorfismos (22.4.5.2) y (22.4.6.2) en condiciones
algo distintas. Sean $A$ un anillo, $\mathfrak m$ un ideal \emph{de cuadrado
nulo} en $A$, $K=A/\mathfrak m$ y $p$ un número primo tal que
$p\cdot1\in\mathfrak m$ (es decir, $K$ es de característica $p$, pero no
necesariamente $A$). Sea, por otra parte, $B$ una $A$-álgebra que sea un $A$-módulo
\emph{fielmente plano} (de modo que $A\subset B$), y supongamos que se tiene
\begin{equation}
\label{0.22.4.8.1-es}
  y^p\in A+pB\subset A+\mathfrak mB
  \tag{22.4.8.1}
\end{equation}
para todo $y\in B$. Si $y^p=x+pz$ con $x\in A$, $z\in B$, la \emph{clase}
mod.~$(A\cap pB)$ de $x$ está bien determinada, y como $A\cap pB=pA$ puesto que
$B$ es un $A$-módulo fielmente plano, se ha definido así una aplicación
canónica $j:B\longrightarrow A/pA$, tal que, para $x\in A$, $j(x)$ sea la
clase de $x^p$ mod.~$pA$. Por otra parte, si $y'\in B$ y
$y'^p=x'+pz'$ con $x'\in A$, $z'\in B$, es inmediato que los elementos
$(y+y')^p-x-x'$ y $(yy')^p-xx'$ pertenecen a $pB$, en virtud de que los
coeficientes binomiales $\binom pi$ son múltiplos de $p$ para
$1\leqslant i\leqslant p-1$. La aplicación $j$ es, pues, un
\emph{homomorfismo de anillos}. Por composición, se deduce de $j$ un
homomorfismo $j':B\xrightarrow{j}A/pA\longrightarrow A/\mathfrak m=K$ y, por
consiguiente, un homomorfismo $j_0:B_{(K)}\longrightarrow K$ tal que $j'$ sea el
compuesto $B\longrightarrow B_{(K)}=B\otimes_AK\xrightarrow{j_0}K$; además,
si $i_0:K\longrightarrow B_{(K)}$ es la aplicación canónica, $i_0$ y $j_0$
verifican las condiciones de (21.3.1) para los anillos \emph{de característica
$p$}, $K$ y $B_{(K)}$. Dicho esto, para definir un $K$-homomorfismo
$\pi'_{B_{(K)}/K}:\Omega^1_{B_{(K)}/K}\otimes_{B_{(K)}}K_{[j_0]}
\longrightarrow\Omega^1_A\otimes_AK$, basta (20.4.8) definir una
$K$-derivación
\begin{equation}
\label{0.22.4.8.2-es}
  D_0:B_{(K)}\longrightarrow\Omega^1_A\otimes_AK
  \tag{22.4.8.2}
\end{equation}
donde el segundo miembro se considera como un $B_{(K)}$-módulo por medio de $j_0$.
Para ello basta definir una $A$-derivación
\begin{equation}
\label{0.22.4.8.3-es}
  D:B\longrightarrow\Omega^1_A\otimes_AK
  \tag{22.4.8.3}
\end{equation}
que \emph{se anule en} $\mathfrak mB$ (considerándose el segundo miembro como
$B$-módulo por medio de $j'$). Ahora bien, si se tiene la relación $y^p=x+pz$ con
$y\in B$, $x\in A$, $z\in B$, el elemento $d_Ax\otimes1_K$ de
$\Omega^1_A\otimes_AK$ solo depende de la clase $j(y)$ mod.~$pA$ de $x$
(por tanto solo de $y$), pues para todo $t\in A$ se tiene
$d_A(pt)=p\,d_A(t)$, por tanto $d_A(pt)\otimes1=0$ en
$\Omega^1_A\otimes_AK=\Omega^1_A/\mathfrak m\Omega^1_A$, puesto que
$p\cdot1\in\mathfrak m$. Se puede tomar, pues,
\[
  D(y)=d_Ax\otimes1_K.
\]
El hecho de que $D$ sea una derivación se deduce fácilmente de que $j$ es un
homomorfismo de anillos; además, para $t\in A$, se tiene
$d_A(t^px)=t^p d_A(x)+pt^{p-1}x d_A(t)$, y de ello se deduce que
$D(ty)=tD(y)$, por tanto $D$ es una $A$-derivación; finalmente, si se tiene además
$t\in\mathfrak m$, se tiene

\oldpage[0\textsubscript{IV-1}]{201}

$D(ty)=tD(y)=0$, puesto que $\Omega^1_A\otimes_AK$ está anulado por $\mathfrak m$.
Se ha definido así el homomorfismo (22.4.5.2) buscado, y es inmediato
verificar que el homomorfismo compuesto
\[
  \Theta_{B_{(K)}/K}\xrightarrow{\pi'_{B_{(K)}/K}}
  \Omega^1_A\otimes_AK\longrightarrow\Omega^1_K
\]
es todavía el homomorfismo canónico $\pi_{B_{(K)}/K}$ (relativo a $i_0$ y
$j_0$). Se tiene, pues, también un homomorfismo canónico análogo a (22.4.5.3).

Si ahora $K$ es un \emph{cuerpo}, $A/pA$ es una $\mathbf F_p$-álgebra,
y se tiene la sucesión exacta (20.5.14)
\[
  0\longrightarrow V/pA\longrightarrow
  \Omega^1_{A/pA}\otimes_{A/pA}K\longrightarrow\Omega^1_K\longrightarrow0
\]
y, por otra parte, se tiene la sucesión exacta (20.5.12.1)
\[
  pA\longrightarrow\Omega^1_A\otimes_A(A/pA)
  \longrightarrow\Omega^1_{A/pA}\longrightarrow0
\]
que muestra, tensorizando con $K$, que los $K$-espacios vectoriales
$\Omega^1_A\otimes_AK$ y
$\Omega^1_{A/pA}\otimes_{A/pA}K$ son isomorfos. El análogo de (22.4.6.2) es,
por tanto, aquí un homomorfismo
\begin{equation}
\label{0.22.4.8.4-es}
  \Xi_{B_{(K)}/K}\longrightarrow V/pA=\mathfrak m/(\mathfrak m^2+pA).
  \tag{22.4.8.4}
\end{equation}

Dicho esto, \emph{el criterio (22.4.7), en el que se sustituyen $V$ y $V'$ por $V/pA$
y $V'/pB$ respectivamente, sigue siendo válido suponiendo solamente que $K$ sea
de característica $p>0$} (es decir, que $p\cdot1\in\mathfrak m$ en
$A$), \emph{y que los $F_i(T_i)$ sean de la forma más general
(22.4.4.1)}: en efecto, $B$ es un $A$-módulo libre (por tanto fielmente plano), y
basta retomar el razonamiento de (22.4.7), sustituyendo (22.4.6.2)
por (22.4.8.4) y $V''$ por $V''/pA''$.
\end{remark}

\subsection{Álgebras geométricamente regulares y álgebras formalmente lisas}
\label{subsection:0.22.5-es}

\begin{proposition}[22.5.1]
\label{0.22.5.1-es}
Sean $A$ un anillo local regular y $A'$ un anillo local que contiene a $A$ y que
es una $A$-álgebra finita, $\mathfrak m$ (resp. $\mathfrak m'$) el ideal maximal
de $A$ (resp. $A'$), y $K=A/\mathfrak m$ (resp. $K'=A'/\mathfrak m'$) su cuerpo
residual. Para que $A'$ sea un anillo regular, es necesario y suficiente que se
tenga
\begin{equation}
\label{0.22.5.1.1-es}
  \operatorname{rg}_K(\mathfrak m/\mathfrak m^2)
  =\operatorname{rg}_{K'}(\mathfrak m'/\mathfrak m'^2).
  \tag{22.5.1.1}
\end{equation}
\end{proposition}

\begin{proof}
En efecto, el primer miembro de (22.5.1.1) es $\dim(A)$ (17.1.1), y como
$A\subset A'$ y $A'$ es una $A$-álgebra finita,
$\dim(A')=\dim(A)$ (16.1.5); la igualdad (22.5.1.1) es, pues, necesaria y
suficiente para que $A'$ sea regular (17.1.1).
\end{proof}

\begin{remark}[22.5.2]
\label{0.22.5.2-es}
\begin{enumerate}
  \item[\rm(i)] Sean $A_1=A/\mathfrak m^2$,
  $A'_1=A'\otimes_AA_1=A'/\mathfrak m^2A'$; como
  $\mathfrak m^2A'\subset\mathfrak m'^2$,
  $\mathfrak m'_1=\mathfrak m'/\mathfrak m^2A'$ es el ideal maximal de $A'_1$
  y $\mathfrak m'_1/\mathfrak m_1'^2$ es un $K'$-espacio vectorial isomorfo
  a $\mathfrak m'/\mathfrak m'^2$; la condición de regularidad de $A'$ solo
  depende, pues, de la estructura de la $A_1$-álgebra $A'_1$, lo que nos
  permitirá a continuación aplicar los resultados preliminares de (22.4) sobre
  las álgebras finitas sobre anillos artinianos.
  \item[\rm(ii)] Se deduce de la demostración de (22.5.1) y de (16.2.6.1)
  que se tiene en todo caso
  \begin{equation}
  \label{0.22.5.2.1-es}
    \operatorname{rg}_K(\mathfrak m/\mathfrak m^2)
    \leqslant\operatorname{rg}_{K'}(\mathfrak m'/\mathfrak m'^2).
    \tag{22.5.2.1}
  \end{equation}

  \oldpage[0\textsubscript{IV-1}]{202}

  \item[\rm(iii)] Supongamos que $A$ y $A'$ verifican las hipótesis generales
  de (22.4.1). Se sabe (19.8.10) que existe un anillo local \emph{regular}
  $B$, de ideal maximal $\mathfrak n$, tal que $A$ es isomorfo a
  $B/\mathfrak n^2$; denotemos por $G_i\in B[T_i]$ un polinomio mónico cuya
  imagen canónica en $A[T_i]$ sea $F_i$ $(1\leqslant i\leqslant r)$;
  entonces, si $B'$ es el anillo cociente de $B[T_1,\ldots,T_r]$ por el ideal
  engendrado por los $G_i$, está claro que $B'$ es un $B$-módulo libre de rango
  $d$ y que $A'=B'\otimes_BA=B'/\mathfrak n^2B'$; como $A'$ se supone un
  anillo local, lo mismo sucede con $B'$, y si $\mathfrak n'$ es el ideal
  maximal de $B'$, $B'/\mathfrak n'$ es isomorfo a $K'$ y
  $\mathfrak n'/\mathfrak n'^2$ es isomorfo a $\mathfrak m'/\mathfrak m'^2$ como
  espacio vectorial sobre $K'$; la desigualdad (22.5.2.1) muestra, pues, que
  se tiene, bajo las hipótesis de (22.4.1),
  \begin{equation}
  \label{0.22.5.2.2-es}
    \operatorname{rg}_K(V)\leqslant\operatorname{rg}_{K'}(\mathfrak m'/\mathfrak m'^2).
    \tag{22.5.2.2}
  \end{equation}
\end{enumerate}
\end{remark}

\begin{corollary}[22.5.3]
\label{0.22.5.3-es}
Sean $A$ un anillo local regular, $\mathfrak m$ su ideal maximal,
$T_i$ $(1\leqslant i\leqslant r)$ indeterminadas y, para cada $i$, sea
\begin{equation}
\label{0.22.5.3.1-es}
  F_i(T_i)=T_i^{d_i}+\sum_{1\leqslant k\leqslant d_i}c_{ik}T_i^{d_i-k}
  \tag{22.5.3.1}
\end{equation}
un polinomio mónico de $A[T_i]$, tal que $d_i\geqslant2$ para todo $i$ y
$c_{ik}\in\mathfrak m$. Sea $A'$ el cociente del anillo de polinomios
$A[T_1,\ldots,T_r]$ por el ideal engendrado por los $r$ polinomios $F_i$. Entonces:
\begin{enumerate}
  \item[\rm(i)] $A'$ es un anillo local, y su cuerpo residual $K'$ es
  isomorfo al cuerpo residual $K$ de $A$.
  \item[\rm(ii)] Para que $A'$ sea un anillo regular, es necesario y suficiente
  que las clases $\xi_i$ mod.~$\mathfrak m^2$ de los elementos $c_{i,d_i}$
  $(1\leqslant i\leqslant r)$ sean linealmente independientes sobre $K$.
\end{enumerate}
\end{corollary}

\begin{proof}
En efecto, con las notaciones de (22.5.2, (i)), si $\overline F_i$ denota el
polinomio de $A_1[T_i]$ obtenido aplicando a los coeficientes de $F_i$
el homomorfismo $A\longrightarrow A_1$, $A'_1$ se define a partir de $A_1$ y
de los polinomios $\overline F_i$ como en (22.4.1), y la conclusión se deduce,
por tanto, de (22.5.1) y de (22.4.3).
\end{proof}

\begin{theorem}[22.5.4]
\label{0.22.5.4-es}
Sean $A$ un anillo local regular y $K$ su cuerpo residual; se supone que
$K$ es de característica $p>0$; sean
\begin{equation}
\label{0.22.5.4.1-es}
  F_i(T_i)=T_i^p+p\sum_{k=1}^{p-1}c_{ik}T_i^{p-k}-f_i
  \tag{22.5.4.1}
\end{equation}
con $c_{ik}\in A$ y $f_i\in A$ $(1\leqslant i\leqslant r)$; sea $B$
el anillo cociente de $A[T_1,\ldots,T_r]$ por el ideal engendrado por los $r$
polinomios $F_i$. Entonces $B$ es un anillo local, y las condiciones siguientes
son equivalentes:
\begin{enumerate}
  \item[a)] $B$ es regular;
  \item[b)] los elementos $d_Af_i\otimes1_K$ $(1\leqslant i\leqslant r)$ son
  linealmente independientes en el $K$-espacio vectorial
  $\Omega^1_A\otimes_AK$;
  \item[c)] si $B'=B/\mathfrak m^2B$ y si se pone
  $\Xi_{B_{(K)}/K}=\Xi_{B'_{(K)}/K}$ (22.4.5), el homomorfismo
  característico $\Xi_{B_{(K)}/K}\longrightarrow\mathfrak m/\mathfrak m^2$
  (22.4.6.2) es inyectivo.
\end{enumerate}
\end{theorem}

\begin{proof}
En efecto, $B'$ se define también a partir de $A_1=A/\mathfrak m^2$ y de los
polinomios $\overline F_i$ obtenidos aplicando a los $F_i$ el homomorfismo
$A\longrightarrow A_1$. Teniendo en cuenta que $\Omega^1_A\otimes_AK$ es
canónicamente isomorfo a $\Omega^1_{A_1}\otimes_{A_1}K$ (20.5.12, (ii)), el
teorema se deduce de (22.5.1) y de (22.4.7.4), utilizando la observación
(22.4.8) cuando $A_1$ no es de característica $p$.
\end{proof}

\oldpage[0\textsubscript{IV-1}]{203}

\begin{env}[22.5.5]
\label{0.22.5.5-es}
Sean $k$ un cuerpo de característica $p>0$, $A$ un anillo local y $k'$ una
extensión \emph{radicial finita} de $k$ tal que $k'^p\subset k$,
$\mathfrak m$ el ideal maximal de $A$ y $K=A/\mathfrak m$ su cuerpo residual;
recordemos que $A'=A\otimes_kk'$ es un anillo local cuyo cuerpo residual es
el mismo que el del anillo local $K\otimes_kk'$ (Bourbaki,
\emph{Alg. conmut.}, cap.~V, §~2, n\textsuperscript{o}~3, lema~4). Observemos que
se tiene $A'_{(K)}=A'\otimes_AK=K\otimes_kk'$, y por consiguiente (21.3.6)
\begin{equation}
\label{0.22.5.5.1-es}
  \Theta_{A'_{(K)}/K}=\Theta_{k'/k}\otimes_kK
  \tag{22.5.5.1}
\end{equation}
salvo isomorfismo canónico. Además, se sabe (21.4.7) que se tiene
$\Xi_{k'/k}=0$, es decir, el homomorfismo canónico (21.3.2.2)
\[
  \pi_{k'/k}:\Theta_{k'/k}\longrightarrow\Omega^1_k
\]
es inyectivo, de modo que se tiene una \emph{inyección canónica}
\begin{equation}
\label{0.22.5.5.2-es}
  \Theta_{A'_{(K)}/K}\longrightarrow\Omega^1_k\otimes_kK
  \tag{22.5.5.2}
\end{equation}
que se explicita de la manera siguiente (teniendo en cuenta (22.5.5.1)): para todo
$x'\in k'$, al elemento
$d_{A'_{(K)}/K}(x'\otimes1_K)\otimes1_K$ de
$\Omega^1_{A'_{(K)}/K}\otimes_{A'_{(K)}}K$ se hace corresponder el elemento
$d_k(x'^p)\otimes1_K$ de $\Omega^1_k\otimes_kK$.
\end{env}

\begin{lemma}[22.5.6]
\label{0.22.5.6-es}
Con las notaciones de (22.5.5):
\begin{enumerate}
  \item[\rm(i)] El núcleo del homomorfismo canónico
  $\pi_{A'_{(K)}/K}$ (21.3.2.2) se identifica por (22.5.5.2) con
  $(\Theta_{k'/k}\otimes_kK)\cap\Upsilon_{K/k}$.
  \item[\rm(ii)] Sean $A_1=A/\mathfrak m^2$,
  $A'_1=A_1\otimes_kk'=A'/\mathfrak m^2A'$; entonces el homomorfismo
  característico $\chi'_{A'_1/A_1}$ (22.4.6.2) se identifica con la restricción del
  homomorfismo característico
  $\chi_{A/k}:\Upsilon_{K/k}\longrightarrow V=\mathfrak m/\mathfrak m^2$
  (20.6.24) a $(\Theta_{k'/k}\otimes_kK)\cap\Upsilon_{K/k}$.
\end{enumerate}
\end{lemma}

\begin{proof}
(i) En virtud de (20.6.9), aplicado sustituyendo $A$ por el cuerpo primo, $B$
por $k$, $C$ por $K$ y $E$ por $A_1=A/\mathfrak m^2$ (extensión de $K$ por
$V=\mathfrak m/\mathfrak m^2$), basta (cf. (22.5.5)) verificar que si
$x'\in k'$ y $z\in A_1$ es un elemento cuya clase mod.
$\mathfrak m/\mathfrak m^2$ es $x'^p\in k$, entonces la imagen de
$d_k(x'^p)\otimes1_K$ por el homomorfismo canónico
$\Omega^1_k\otimes_kK\longrightarrow\Omega^1_{A_1}\otimes_{A_1}K$ es
$d_{A_1}z\otimes1_K$, lo que se deduce de las definiciones.

(ii) Como $A'_{(K)}=K\otimes_kk'=(A'_1)_{(K)}$ (en virtud de la relación
$K=A_1/(\mathfrak m/\mathfrak m^2)$), $\Theta_{A'_{(K)}/K}$ se identifica con
$\Theta_{(A'_1)_{(K)}/K}$ y la verificación resulta del mismo cálculo que en
(i) y de la definición (22.4.5.2) de $\pi'_{(A'_1)_{(K)}/K}$.
\end{proof}

\begin{proposition}[22.5.7]
\label{0.22.5.7-es}
Sean $k$ un cuerpo de característica $p>0$ y $A$ una $k$-álgebra que es un
anillo local regular, $\mathfrak m$ su ideal maximal y $K=A/\mathfrak m$ su
cuerpo residual; sea $V=\mathfrak m/\mathfrak m^2$, considerado como
$K$-espacio vectorial. Sean $k'$ una extensión finita de $k$ tal que
$k'^p\subset k$, y $A'=A\otimes_kk'$; para que el anillo local $A'$ sea
regular, es necesario y suficiente que la restricción del homomorfismo
característico $\chi_{A/k}:\Upsilon_{K/k}\longrightarrow V$ (20.6.24) a
$(\Theta_{k'/k}\otimes_kK)\cap\Upsilon_{K/k}$ sea inyectiva.
\end{proposition}

\begin{proof}
Se trata de demostrar que la condición del enunciado equivale a
(22.5.1.1), denotando por $\mathfrak m'$ el ideal maximal de $A'$ y por
$K'=A'/\mathfrak m'$ su cuerpo residual. Con las notaciones de la observación
(22.5.2, (i)), se tiene $A'_1=A_1\otimes_kk'$. Ahora bien, sea
$(a_i)_{1\leqslant i\leqslant r}$ una $p$-base de $k'$ sobre $k$, de modo que
$a_i^p=b_i\in k$, y $k'$ es isomorfo al cociente del anillo de

\oldpage[0\textsubscript{IV-1}]{204}

polinomios $k[T_1,\ldots,T_r]$ por el ideal engendrado por los polinomios
$F_i=T_i^p-b_i$ $(1\leqslant i\leqslant r)$; se deduce que $A'_1$ es
isomorfo al cociente del anillo de polinomios $A_1[T_1,\ldots,T_r]$ por
el ideal engendrado por los $F_i$. Se puede aplicar entonces el criterio (22.4.7), y
la condición del enunciado equivale, por tanto, a que $\chi'_{A'_1/A_1}$ sea
inyectivo; se concluye con ayuda de (22.5.6).
\end{proof}

Podemos ahora demostrar el recíproco de (19.6.5):

\begin{theorem}[22.5.8]
\label{0.22.5.8-es}
Sean $k$ un cuerpo, $p$ su exponente característico y $A$ una $k$-álgebra
local noetheriana. Las condiciones siguientes son equivalentes:
\begin{enumerate}
  \item[a)] $A$ es una $k$-álgebra formalmente lisa (para su topología
  preádica).
  \item[b)] $A$ es geométricamente regular sobre $k$.
  \item[b')] Para toda extensión finita $k'$ de $k$ tal que $k'^p\subset k$,
  $A'=A\otimes_kk'$ es un anillo regular.
\end{enumerate}
\end{theorem}

\begin{proof}
Teniendo en cuenta (19.6.6), se puede reducir al caso $p>1$.

El hecho de que a) implique b) no es otro que (19.6.5), y b) implica
trivialmente b'). Demostremos que b') implica que se satisfacen las condiciones de
(22.2.2); esto es evidente para la primera (17.1.1), puesto que b') implica
en primer lugar que $A$ es regular. Por otra parte, sea $(x_\alpha)_{\alpha\in I}$ una
$p$-base de $k^{1/p}$ sobre $k$; se sabe que los elementos $d_k(x_\alpha^p)$
forman una base del $k$-espacio vectorial $\Omega^1_k$ (21.4.5), por lo que los
elementos $d_k(x_\alpha^p)\otimes1_K$ forman una base del $K$-espacio vectorial
$\Omega^1_k\otimes_kK$. Se concluye (cf. (22.5.5)) que cuando $k'$ recorre
el conjunto de las subextensiones de $k^{1/p}$, finitas sobre $k$, la familia de
subespacios $\Theta_{k'/k}\otimes_kK$ de $\Omega^1_k\otimes_kK$ es filtrante
creciente y tiene por unión $\Omega^1_k\otimes_kK$. Se deduce entonces de
(22.5.7) que la condición b') implica que $\chi_{A/k}$ es inyectivo, lo que
no es sino la condición (ii) de (22.2.2).
\end{proof}

\begin{corollary}[22.5.9]
\label{0.22.5.9-es}
Sean $k$ un cuerpo y $A$ una $k$-álgebra local noetheriana formalmente
lisa. Entonces, para todo ideal primo $\mathfrak p$ de $A$, el anillo local
$A_{\mathfrak p}$ es una $k$-álgebra formalmente lisa (para la topología
$\mathfrak p$-preádica).
\end{corollary}

\begin{proof}
En efecto, con las notaciones de (22.5.8, b')), basta demostrar que
el anillo local $A_{\mathfrak p}\otimes_kk'$ es regular; pero como es un
anillo de fracciones del anillo local $A\otimes_kk'$ y este último es
regular por hipótesis, lo mismo sucede con $A_{\mathfrak p}\otimes_kk'$
(17.3.6).
\end{proof}

\begin{remark}[22.5.10]
\label{0.22.5.10-es}
\begin{enumerate}
  \item[\rm(i)] La conclusión de (22.5.8) falla cuando, en la
  condición b'), uno se limita a las extensiones \emph{monógenas} $k'=k(x)$ de $k$
  (con $x^p\in k$). Puede suceder, en efecto, que para todas estas extensiones $k'$,
  se tenga
  $\Upsilon_{K/k}\cap(\Theta_{k'/k}\otimes_kK)=0$, aunque
  $\Upsilon_{K/k}\neq0$; esto significa que para todo $x\in k^{1/p}$ tal que
  $x\notin k$, debe tenerse
  $d_k(x^p)\otimes1_K\notin\Upsilon_{K/k}$, o también $d_K(x^p)\neq0$,
  es decir, $x^p\notin K^p$; en otras palabras, debe tenerse
  $K^p\cap k=k^p$, aunque $K$ sea una extensión \emph{inseparable} de $k$.
  Ahora bien, se construyen fácilmente ejemplos de tales extensiones: partamos de un
  cuerpo perfecto $k_0$ de característica $p>0$, sean $X$, $Y$, $Z$ tres
  indeterminadas y pongamos $k=k_0(X,Y)$ y
  $K=k(X^{1/p}+ZY^{1/p})$; se verifica fácilmente que $K$ satisface las condiciones
  precedentes, por lo que $\Upsilon_{K/k}\neq0$. Apliquemos ahora el lema
  (22.2.5.2) tomando $V=K$ y como $h$ el homomorfismo nulo; $A$ es entonces
  un anillo de valoración discreta que tiene a $K$ por cuerpo residual, con
  $\chi_{A/k}=0$, y no es, por tanto, una $k$-álgebra formalmente lisa por
  (22.2.2); sin embargo, $A\otimes_kk'$ es un anillo regular para toda
  extensión \emph{monógena} $k'\subset k^{1/p}$ de $k$.

  \oldpage[0\textsubscript{IV-1}]{205}

  \item[\rm(ii)] El corolario (22.5.9) lleva a considerar la cuestión
  siguiente: si $A$ es una $k$-álgebra noetheriana, ¿bajo qué condiciones
  el conjunto de los ideales primos $\mathfrak p\in\operatorname{Spec}(A)$ tales que
  $A_{\mathfrak p}$ sea una $k$-álgebra formalmente lisa es
  \emph{abierto}? Abordaremos más adelante algunos casos particulares.
\end{enumerate}
\end{remark}
\subsection{Criterio jacobiano de Zariski}
\label{subsection:0.22.6-es}

\begin{theorem}[22.6.1]
\label{0.22.6.1-es}
\textnormal{(criterio jacobiano de lisitud formal).}
Sean $A$, $B$ dos anillos topológicos,
$u:A\longrightarrow B$ un homomorfismo continuo que hace de $B$ una
$A$-álgebra formalmente lisa, $\mathfrak J$ un ideal de $B$ (no
necesariamente cerrado), $C=B/\mathfrak J$ el $A$-álgebra topológica cociente.
Para que $C$ sea una $A$-álgebra formalmente lisa, es necesario y suficiente que
el homomorfismo canónico (cf. (20.5.11.2))
\begin{equation}
\label{0.22.6.1.1-es}
  \delta_{C/B/A}:\mathfrak J/\mathfrak J^2
  \longrightarrow\Omega^1_{B/A}\otimes_BC
  \tag{22.6.1.1}
\end{equation}
sea \emph{formalmente invertible por la izquierda} (cf. (19.1.5)).
\end{theorem}

\begin{proof}
En efecto, decir que $B$ (resp. $C$) es una $A$-álgebra formalmente lisa
significa (19.4.4) que se tiene
$\operatorname{Exalcotop}_A(B,L)=0$ (resp.
$\operatorname{Exalcotop}_A(C,L)=0$) para todo $B$-módulo (resp. $C$-módulo)
discreto $L$ anulado por un ideal abierto de $B$ (resp. $C$). Como, por hipótesis,
$\operatorname{Exalcotop}_A(B,L)=0$ para todo $C$-módulo discreto $L$ anulado
por un ideal abierto de $C$, decir que $C$ es una $A$-álgebra formalmente lisa
significa, pues, que el homomorfismo canónico
$\operatorname{Exalcotop}_A(C,L)\longrightarrow
\operatorname{Exalcotop}_A(B,L)$ es inyectivo, y por tanto el teorema resulta de
(20.7.8).
\end{proof}

Recordemos (19.5.3) que, cuando $B$ y $C$ son $A$-álgebras formalmente
lisas, $\mathfrak J/\mathfrak J^2$ es un $C$-módulo topológico
\emph{formalmente proyectivo}; además, los homomorfismos canónicos
$\varphi_n:\mathbf S_C^n(\mathfrak J/\mathfrak J^2)
\longrightarrow\mathfrak J^n/\mathfrak J^{n+1}$ son
\emph{bimorfismos formales} cuando se supone que $B$ es un anillo
\emph{preadmisible}.
\begin{corollary}[22.6.2]
\label{0.22.6.2-es}
Siendo las hipótesis y notaciones las de (22.6.1), supongamos además que
en $B$ el cuadrado de todo ideal abierto es abierto, y que sobre $\mathfrak J$ la
topología inducida por la de $B$ es idéntica a la topología deducida de
la de $B$ (19.0) (se observará que estas dos condiciones se verifican cuando
$B$ es noetheriano y su topología es preádica ($0_{\mathrm I}$, 7.3.2), o si la
topología de $B$ es la topología $\mathfrak J$-preádica). Sea
$(\mathfrak b_\lambda)$ un sistema fundamental de ideales abiertos de $B$ y
pongamos $B_\lambda=B/\mathfrak b_\lambda$ para todo $\lambda$. Entonces:
\begin{enumerate}
  \item[\rm(i)] Las condiciones siguientes son equivalentes:
  \begin{enumerate}
    \item[a)] $C$ es una $A$-álgebra formalmente lisa.
    \item[b)] Para todo $\lambda$, el homomorfismo de $B_\lambda$-módulos
    \begin{equation}
    \label{0.22.6.2.1-es}
      (\mathfrak J/\mathfrak J^2)\otimes_BB_\lambda
      \longrightarrow\Omega^1_{B/A}\otimes_BB_\lambda
      \tag{22.6.2.1}
    \end{equation}
    deducido de $\delta_{C/B/A}$ por tensorización, es invertible por la izquierda
    (en otras palabras, un isomorfismo sobre un factor directo de
    $\Omega^1_{B/A}\otimes_BB_\lambda$).
  \end{enumerate}
  \item[\rm(ii)] Supongamos además que el anillo topológico $C$ es
  preadmisible ($0_{\mathrm I}$, 7.1.2) y

  \oldpage[0\textsubscript{IV-1}]{206}

  sea $\mathfrak L$ un ideal de definición de $C$; entonces las condiciones a)
  y b) de (i) también equivalen a:
  \begin{enumerate}
    \item[c)] El homomorfismo de $(C/\mathfrak L)$-módulos
    \[
      (\mathfrak J/\mathfrak J^2)\otimes_C(C/\mathfrak L)
      \longrightarrow\Omega^1_{B/A}\otimes_B(C/\mathfrak L)
    \]
    es invertible por la izquierda.
  \end{enumerate}
\end{enumerate}
\end{corollary}

\begin{proof}
Las hipótesis implican que sobre $\Omega^1_{B/A}$ la topología es deducida de
la de $B$ (20.4.5), por lo que (i) resulta de (22.6.1) y de (19.1.7). Por otra
parte, recordemos (20.4.9) que, puesto que $B$ es una $A$-álgebra formalmente
lisa, $\Omega^1_{B/A}$ es un $B$-módulo formalmente proyectivo, luego
$\Omega^1_{B/A}\otimes_BB_\lambda$ es un $B_\lambda$-módulo proyectivo para
todo $\lambda$ (19.2.4); la equivalencia de c) con a) y b) cuando $C$ es
preadmisible resulta entonces de (19.1.9).
\end{proof}

En particular:

\begin{corollary}[22.6.3]
\label{0.22.6.3-es}
Sean $A$ un anillo, $B$ una $A$-álgebra, $\mathfrak J$ un ideal de $B$,
$C=B/\mathfrak J$; se supone que $A$, $B$, $C$ están provistos de las topologías discretas
y que $B$ es una $A$-álgebra formalmente lisa. Para que $C$ sea una
$A$-álgebra formalmente lisa, es necesario y suficiente que el homomorfismo
canónico (22.6.1.1) sea invertible por la izquierda.
\end{corollary}

Se observará que, puesto que toda $A$-álgebra $C$ es cociente de un álgebra de
polinomios $A[T_\alpha]_{\alpha\in I}$, y que esta última es formalmente
lisa (19.3.3), el criterio (22.6.3) permite en principio reconocer si $C$
es una $A$-álgebra formalmente lisa.

\begin{proposition}[22.6.4]
\label{0.22.6.4-es}
Sean $A$ un anillo, $B$ una $A$-álgebra formalmente lisa (para las
topologías discretas), $\mathfrak J$ un ideal de $B$, $C=B/\mathfrak J$; se
supone que $\mathfrak J/\mathfrak J^2$ es un $C$-módulo de tipo finito. Sean
$\mathfrak p$ un ideal primo de $C$, $k(\mathfrak p)$ el cuerpo residual de
$C_{\mathfrak p}$, $\mathfrak p'$ la imagen inversa de $\mathfrak p$ en
$B$, $\mathfrak q$ el ideal primo de $A$ imagen inversa de $\mathfrak p'$.
Las condiciones siguientes son equivalentes:
\begin{enumerate}
  \item[a)] $C_{\mathfrak p}$ es una $A_{\mathfrak q}$-álgebra formalmente
  lisa (para las topologías discretas).
  \item[a')] $C_{\mathfrak p}$ es una $A$-álgebra formalmente lisa (para
  las topologías discretas).
  \item[b)] El homomorfismo canónico
  \begin{equation}
  \label{0.22.6.4.1-es}
    (\mathfrak J/\mathfrak J^2)\otimes_Ck(\mathfrak p)
    \longrightarrow\Omega^1_{B/A}\otimes_Bk(\mathfrak p)
    \tag{22.6.4.1}
  \end{equation}
  es inyectivo.
  \item[c)] Existe $f\in C-\mathfrak p$ tal que $C_f$ es una
  $A$-álgebra formalmente lisa.
\end{enumerate}
Si además $B$ es noetheriano, las condiciones anteriores también son
equivalentes a:
\begin{enumerate}
  \item[d)] $C_{\mathfrak p}$ es una $A_{\mathfrak q}$-álgebra formalmente
  lisa para la topología $\mathfrak p$-preádica sobre $C_{\mathfrak p}$ y la
  topología discreta (o la topología $\mathfrak q$-preádica) sobre
  $A_{\mathfrak q}$.
\end{enumerate}
\end{proposition}

\begin{proof}
Como $A_{\mathfrak q}$ es una $A$-álgebra formalmente lisa (19.3.5, (iv)
y (ii)), ya se sabe que a) y a') son equivalentes (19.3.5, (ii)). Se tiene
$C_{\mathfrak p}=B_{\mathfrak p'}/\mathfrak J B_{\mathfrak p'}$, y
$B_{\mathfrak p'}$ es una $A$-álgebra formalmente lisa para la topología
discreta (19.3.5, (iv)); además,
$\Omega^1_{B_{\mathfrak p'}/A}=(\Omega^1_{B/A})_{\mathfrak p'}$ (20.5.9).
La equivalencia de a') y b) resulta entonces de aplicar (22.6.3) a
$B_{\mathfrak p'}$ y $C_{\mathfrak p}$, teniendo en cuenta que
$\Omega^1_{B/A}$ es un $B$-módulo proyectivo (20.4.9) y (19.2.11), y que
$\mathfrak J/\mathfrak J^2$ es un $B$-módulo de tipo finito, y utilizando
(19.1.12). La aplicación de (19.1.12) prueba también la equivalencia de b) y c).
Por último, observemos que, puesto que $B_{\mathfrak p'}$ es una
$A_{\mathfrak q}$-álgebra formalmente lisa para las topologías

\oldpage[0\textsubscript{IV-1}]{207}

discretas (19.3.5, (iv)), $B_{\mathfrak p'}$, provisto de la topología
$\mathfrak p'$-preádica, sigue siendo una $A_{\mathfrak q}$-álgebra
formalmente lisa cuando se toma sobre $A_{\mathfrak q}$ la topología
discreta o la topología $\mathfrak q$-preádica (19.3.8). Para demostrar
la equivalencia de b) y d) cuando $B$ es noetheriano, apliquemos (22.6.2) al
anillo $A_{\mathfrak q}$ discreto (o $\mathfrak q$-preádico) y al anillo
$B_{\mathfrak p'}$, provisto de la topología $\mathfrak p'$-preádica; como
$C_{\mathfrak p}$ es preadmisible y $\mathfrak pC_{\mathfrak p}$ es
un ideal de definición para su topología, se puede invocar la equivalencia de
c) y a) en (22.6.2).
\end{proof}

\begin{corollary}[22.6.5]
\label{0.22.6.5-es}
Bajo las hipótesis generales de (22.6.4), el conjunto de los
$\mathfrak p\in\operatorname{Spec}(C)$ tales que $C_{\mathfrak p}$ sea una $A$-álgebra
formalmente lisa (o una $A_{\mathfrak q}$-álgebra formalmente lisa, designando
por $\mathfrak q$ la imagen inversa de $\mathfrak p$ en $A$)
para las topologías discretas, es abierto en $\operatorname{Spec}(C)$.
\end{corollary}

Esto resulta de la forma c) de (22.6.4). Se observará que, cuando $B$ es
noetheriano, se pueden sustituir las topologías discretas por las topologías
preádicas.

\begin{corollary}[22.6.6]
\label{0.22.6.6-es}
Bajo las hipótesis generales de (22.6.4), las condiciones siguientes son
equivalentes:
\begin{enumerate}
  \item[a)] $C$ es una $A$-álgebra formalmente lisa (para las topologías
  discretas).
  \item[b)] Para todo $\mathfrak p\in\operatorname{Spec}(C)$ (o solamente para todo ideal
  máximo $\mathfrak p$ de $C$), $C_{\mathfrak p}$ es una $A$-álgebra
  formalmente lisa (o una $A_{\mathfrak q}$-álgebra formalmente lisa, designando
  por $\mathfrak q$ la imagen inversa de $\mathfrak p$ en $A$)
  para las topologías discretas.
\end{enumerate}
Además, cuando $B$ es noetheriano, se pueden sustituir en b) las topologías
discretas por las topologías preádicas sobre $A_{\mathfrak q}$ y
$C_{\mathfrak p}$.
\end{corollary}

\begin{proof}
La última afirmación resulta de (22.6.4). En virtud de (22.6.3), la condición
a) equivale a que el homomorfismo (22.6.1.1) sea invertible por la izquierda;
la equivalencia de a) y b) resulta, pues, de (19.1.14), teniendo en cuenta que
$\Omega^1_{B/A}$ es un $B$-módulo proyectivo y $\mathfrak J/\mathfrak J^2$
un $C$-módulo de tipo finito.
\end{proof}

\begin{proposition}[22.6.7]
\label{0.22.6.7-es}
Sean $k_0$ un cuerpo, $k$ una extensión separable de $k_0$, $C$ una
$k$-álgebra de tipo finito.
\begin{enumerate}
  \item[\rm(i)] Sea $\mathfrak p$ un ideal primo de $C$. Las condiciones
  siguientes son equivalentes:
  \begin{enumerate}
    \item[a)] $C_{\mathfrak p}$ es una $k_0$-álgebra formalmente lisa para
    la topología discreta.
    \item[b)] $C_{\mathfrak p}$ es una $k_0$-álgebra formalmente lisa para
    la topología $\mathfrak p$-preádica.
    \item[c)] Existe $f\in C-\mathfrak p$ tal que $C_f$ es una
    $k_0$-álgebra formalmente lisa para la topología discreta.
    \item[d)] $C_{\mathfrak p}$ es un anillo geométricamente regular
    (19.6.5) sobre $k_0$.
  \end{enumerate}
  Además, el conjunto de los $\mathfrak p\in\operatorname{Spec}(C)$ que poseen cualquiera de
  estas propiedades es abierto en $\operatorname{Spec}(C)$.
  \item[\rm(ii)] Las condiciones siguientes son equivalentes:
  \begin{enumerate}
    \item[a)] $C$ es una $k_0$-álgebra formalmente lisa para la topología
    discreta.
    \item[b)] Todo $\mathfrak p\in\operatorname{Spec}(C)$ verifica las condiciones
    equivalentes de (i).
    \item[c)] Todo ideal máximo $\mathfrak m$ de $C$ verifica las condiciones
    equivalentes de (i).
    \item[d)] Para toda extensión $k'$ de $k_0$, $C\otimes_{k_0}k'$ es un
    anillo regular (17.3.6); se dice también que $C$ es un anillo
    geométricamente regular sobre $k_0$.
  \end{enumerate}

  \oldpage[0\textsubscript{IV-1}]{208}

  \item[\rm(iii)] Sean $B=k[T_1,\ldots,T_r]$ un álgebra de polinomios,
  $\mathfrak J$ un ideal de $B$ tales que $C$ sea isomorfa a
  $B/\mathfrak J$. Sea $\mathfrak q$ un ideal primo de $B$ que contiene
  $\mathfrak J$ y sea $\mathfrak p=\mathfrak q/\mathfrak J$. Las condiciones
  de (i) son entonces también equivalentes a la siguiente («criterio jacobiano de
  Zariski»):
  \begin{enumerate}
    \item[e)] Existe un sistema de elementos
    $(f_i)_{1\leq i\leq m}$ de $\mathfrak J$, que engendra el ideal
    $\mathfrak J_{\mathfrak q}$ en $B_{\mathfrak q}$, y de
    $k_0$-derivaciones $D_j$ de $B$ en sí mismo $(1\leq j\leq m)$ tales que
    $\det(D_j(f_i))\notin\mathfrak q$.
  \end{enumerate}
\end{enumerate}
\end{proposition}

\begin{proof}
Se sabe que $C$ es siempre isomorfa a una $k$-álgebra de la forma
$B/\mathfrak J$. Como $B$ es una $k$-álgebra formalmente lisa (19.3.3) y
$k$, al ser separable sobre $k_0$, es una $k_0$-álgebra formalmente lisa
(19.6.1), $B$ es una $k_0$-álgebra formalmente lisa (19.3.5, (ii)).
La equivalencia de las condiciones a), b), c) de (i) resulta, pues, de (22.6.4), y
la de a), b), c) en (ii) resulta de (22.6.6); la equivalencia de a) y d)
en (i) resulta de (22.5.8); como todo localizado de $C\otimes_{k_0}k'$ es
también un localizado de $C_{\mathfrak p}\otimes_{k_0}k'$ para un ideal primo
$\mathfrak p\in\operatorname{Spec}(C)$ conveniente, la equivalencia de d) y b) en (ii)
resulta de la equivalencia de d) y a) en (i). Finalmente, como
$\Omega^1_{B/k_0}$ es un $B$-módulo proyectivo, la equivalencia del criterio de
Zariski e) y de las demás condiciones de (i) se deduce de (19.1.12) y de
(22.6.3), teniendo en cuenta (20.4.8).
\end{proof}

Zariski se interesaba, de hecho, por un criterio diferencial de regularidad para los
anillos locales $C_{\mathfrak p}$. Como es equivalente decir que un anillo
local noetheriano que contiene un cuerpo es regular o es formalmente liso en
tanto que álgebra sobre su subcuerpo primo (19.6.4), se obtiene inmediatamente tal
criterio sustituyendo $k_0$ por el subcuerpo primo de $k$ en (22.6.7); en
particular, se obtiene el resultado siguiente, que volveremos a encontrar más adelante
(IV, 6.12.5) como caso particular de resultados más generales de Nagata:

\begin{corollary}[22.6.8]
\label{0.22.6.8-es}
\textnormal{(Zariski).}
Sea $C$ un álgebra de tipo finito sobre un cuerpo. Entonces el conjunto de los
$\mathfrak p\in\operatorname{Spec}(C)$ tales que $C_{\mathfrak p}$ sea un anillo local
regular es abierto en $\operatorname{Spec}(C)$.
\end{corollary}

\begin{remark}[22.6.9]
\label{0.22.6.9-es}
\begin{enumerate}
  \item[\rm(i)] Los enunciados de (22.6.7) siguen siendo válidos si, en lugar de
  suponer $k$ separable sobre $k_0$, se supone solamente que es de
  \emph{multiplicidad radicial finita}, es decir (19.6.6), que existe
  una extensión radicial finita $k'_0$ de $k_0$ tal que el cuerpo residual del
  anillo artiniano local $k\otimes_{k_0}k'_0$ sea una extensión separable
  $k'$ de $k'_0$. Existe entonces un $k'_0$-monomorfismo
  $k'\longrightarrow k\otimes_{k_0}k'_0$ que, compuesto con el homomorfismo
  canónico $k\otimes_{k_0}k'_0\longrightarrow k'$, da la identidad, puesto que
  $k'$ es una $k'_0$-álgebra formalmente lisa (19.6.1); se concluye que
  $C\otimes_{k_0}k'_0$ está provista de una estructura de $k'$-álgebra de tipo finito,
  y, como $k'$ es separable sobre $k'_0$, se puede aplicar (22.6.7) a esta
  $k'$-álgebra; se concluye nuestra afirmación aplicando (19.4.7). Como
  no habremos de utilizar esta generalización, dejamos los detalles al
  lector. Ignoramos, en cambio, si los resultados de (22.6.7) son válidos sin
  \emph{ninguna} hipótesis sobre la extensión $k$ de $k_0$.
  \item[\rm(ii)] Bajo las hipótesis generales de (22.6.7), sea $(k_\alpha)$
  una familia filtrante decreciente de subcuerpos de $k$ que contienen $k_0$ y
  tal que
  $\bigcap_\alpha k_\alpha(k^p)=k_0(k^p)$ (donde $p$ es el exponente
  característico de $k_0$). Utilizando un método de cómputo de dimensiones
  debido a Nagata, se puede demostrar que, en el criterio jacobiano (22.6.7, (iii)),
  se puede uno limitar a derivaciones $D_j$ de $B$ que sean
  $k_\alpha$-derivaciones para un $\alpha$ conveniente.

  \oldpage[0\textsubscript{IV-1}]{209}

  El interés de este resultado es que siempre existen tales familias
  $(k_\alpha)$ para las cuales $[k:k_\alpha]$ es finito (21.8.6). No
  demostraremos este refinamiento del criterio de Zariski, que no habremos de
  utilizar. En (22.7), vamos a dar, para los anillos locales completos,
  una variante (debida también a Nagata) del criterio de Zariski, que se demuestra
  esencialmente por el mismo método (con dificultades técnicas algo
  mayores).
\end{enumerate}
\end{remark}

\subsection{El criterio jacobiano de Nagata}
\label{subsection:0.22.7-es}

\begin{env}[22.7.1]
\label{0.22.7.1-es}
El criterio jacobiano de Nagata es el análogo del criterio jacobiano de Zariski,
pero para anillos cocientes de anillos de \emph{series formales} sobre un
cuerpo. Daremos, como Nagata [31], dos versiones, presentadas aquí
como criterios de lisitud formal.
\end{env}

\begin{proposition}[22.7.2]
\label{0.22.7.2-es}
Sean $k$ un cuerpo, $A=k[[T_1,\ldots,T_r]]$ un anillo de series formales sobre
$k$, provisto de su estructura usual de $k$-álgebra, $\mathfrak q$ un ideal de
$A$, $B=A/\mathfrak q$. Sea $\mathfrak p$ un ideal primo de $A$ que contiene
$\mathfrak q$. Se supone que existe una sub-$k$-álgebra $C'$ de
$C=A/\mathfrak p$, isomorfa a un álgebra de series formales
$k[[X_1,\ldots,X_s]]$, tal que $C$ sea finita sobre $C'$ y que el cuerpo de
fracciones $L$ de $C$ sea una extensión separable de
$L'=k((X_1,\ldots,X_s))$ (hipótesis siempre satisfecha si $k$ es de
característica 0). Entonces las condiciones siguientes son equivalentes:
\begin{enumerate}
  \item[a)] $B_{\mathfrak p}=A_{\mathfrak p}/\mathfrak qA_{\mathfrak p}$ es
  una $k$-álgebra formalmente lisa para la topología
  $\mathfrak p$-preádica.
  \item[b)] Existen $k$-derivaciones $D_i$ $(1\leq i\leq m)$ de $A$ en
  sí mismo, y elementos $f_i$ $(1\leq i\leq m)$ de $\mathfrak q$, tales
  que las imágenes de los $f_i$ en $\mathfrak qA_{\mathfrak p}$ engendran este
  ideal de $A_{\mathfrak p}$, y se tiene
  $\det(D_i(f_j))\notin\mathfrak p$.
  \item[c)] $B_{\mathfrak p}$ es un anillo regular.
\end{enumerate}
\end{proposition}

\begin{proof}
El cuerpo residual de $B_{\mathfrak p}$ es $k(\mathfrak p)$; como
$k((X_1,\ldots,X_s))$ es separable sobre $k$ (21.9.6.4), $k(\mathfrak p)$ es
separable sobre $k$ por hipótesis, y ya se sabe que en estas condiciones las
propiedades a) y c) son equivalentes (19.6.4). Por otra parte,
$A_{\mathfrak p}$ es un anillo regular ((17.1.4) y (17.3.2)), y como su
cuerpo residual $k(\mathfrak p)$ es separable sobre $k$, $A_{\mathfrak p}$ es
formalmente lisa sobre $k$ para la topología $\mathfrak p$-preádica
(19.6.4). Resulta, pues, de (22.6.2, (ii)) que la condición a) equivale a que
el homomorfismo de $L$-espacios vectoriales
\begin{equation}
\label{0.22.7.2.1-es}
  j_0:(\mathfrak q/\mathfrak q^2)\otimes_AL
  \longrightarrow\Omega^1_{A_{\mathfrak p}/k}\otimes_{A_{\mathfrak p}}L
  =\widehat\Omega^1_{A/k}\otimes_AL
  \tag{22.7.2.1}
\end{equation}
sea inyectivo.

Consideremos, por otra parte, el homomorfismo compuesto
\begin{equation}
\label{0.22.7.2.2-es}
  j:(\mathfrak q/\mathfrak q^2)\otimes_AL
  \xrightarrow{j_0}\Omega^1_{A/k}\otimes_AL
  \longrightarrow\widehat\Omega^1_{A/k}\otimes_AL
  \tag{22.7.2.2}
\end{equation}
y demostremos que la condición b) equivale a decir que $j$ es inyectivo. Observemos, en
efecto, que
\[
  (\mathfrak q/\mathfrak q^2)\otimes_AL
  =(\mathfrak qA_{\mathfrak p}/\mathfrak q^2A_{\mathfrak p})
  \otimes_{A_{\mathfrak p}}L~;
\]
la condición b) significa que, si $F_i=j(f'_i\otimes1)$, donde $f'_i$ es la imagen
de $f_i$ en
$\mathfrak qA_{\mathfrak p}/\mathfrak q^2A_{\mathfrak p}$, la matriz
$(\langle F_i,D_j\otimes1\rangle)$ es invertible; por tanto, esto implica que los
$F_i$ son linealmente independientes, y \emph{a fortiori} que ocurre lo
mismo con los $f'_i\otimes1$; pero como estos últimos engendran
$(\mathfrak qA_{\mathfrak p}/\mathfrak q^2A_{\mathfrak p})
\otimes_{A_{\mathfrak p}}L$, se concluye que $j$ es inyectivo. Recíprocamente,
supongamos $j$ inyectivo, y tomemos los $f_i$ de modo que los $f'_i\otimes1$
formen

\oldpage[0\textsubscript{IV-1}]{210}

una base de
$(\mathfrak qA_{\mathfrak p}/\mathfrak q^2A_{\mathfrak p})
\otimes_{A_{\mathfrak p}}L$; entonces los $F_i$ forman una base de la imagen de
$j$; pero se sabe (21.9.3) que $\widehat\Omega^1_{A/k}$ es un $A$-módulo
libre de rango $r$ y que las $k$-derivaciones de $A$ en sí mismo engendran su
dual; el hecho de que los $F_i$ sean linealmente independientes implica, pues,
la existencia de $k$-derivaciones $D_j$ de $A$ en sí mismo tales que la matriz
$(\langle F_i,D_j\otimes1\rangle)$ sea invertible, es decir, la condición
b).

Estas observaciones muestran ya que b) implica a). Para probar recíprocamente que
la condición a) implica que $j$ es inyectivo, consideremos el diagrama
conmutativo
\begin{equation}
\label{0.22.7.2.3-es}
\xymatrix{
  (\mathfrak q/\mathfrak q^2)\otimes_AL
    \ar[r]^j \ar[d]_\alpha &
  \widehat\Omega^1_{A/k}\otimes_AL \ar@{=}[d] \\
  (\mathfrak p/\mathfrak p^2)\otimes_AL
    \ar[r]_i &
  \widehat\Omega^1_{A/k}\otimes_AL
}
\tag{22.7.2.3}
\end{equation}
y observemos el lema siguiente:
\end{proof}

\begin{lemma}[22.7.2.4]
\label{0.22.7.2.4-es}
Sean $R$ un anillo local regular, $\mathfrak m$ su ideal máximo,
$K=R/\mathfrak m$ su cuerpo residual. Si $\mathfrak n$ es un ideal de $R$ tal
que $R/\mathfrak n$ sea regular, entonces el homomorfismo canónico
\begin{equation}
\label{0.22.7.2.5-es}
  \alpha:(\mathfrak n/\mathfrak n^2)\otimes_RK
  \longrightarrow(\mathfrak m/\mathfrak m^2)\otimes_RK
  \tag{22.7.2.5}
\end{equation}
es inyectivo.
\end{lemma}

\begin{proof}
En efecto, se sabe (17.1.6) que $\mathfrak n$ está engendrado por una sucesión
$(x_i)_{1\leq i\leq t}$ que forma parte de un sistema regular de parámetros
$(x_i)_{1\leq i\leq n}$ de $R$; las imágenes de los $x_i$ $(1\leq i\leq t)$
forman una base de $(\mathfrak n/\mathfrak n^2)\otimes_RK$, y sus imágenes
por $\alpha$ forman parte de la base de
$(\mathfrak m/\mathfrak m^2)\otimes_RK$ constituida por las imágenes de los $x_i$ para
$1\leq i\leq n$, de donde resulta el lema.
\end{proof}

\begin{proof}[Continuación de la demostración de la proposición~{\rm(22.7.2)}]
Este lema y el diagrama (22.7.2.3) reducen, pues, (en virtud de la hipótesis a))
la cuestión a probar que $i$ es inyectivo. Ahora bien, en la sucesión
\begin{equation}
\label{0.22.7.2.6-es}
  \mathfrak p/\mathfrak p^2\xrightarrow{h}
  \widehat\Omega^1_{A/k}\widehat\otimes_A(A/\mathfrak p)
  \xrightarrow{g}\widehat\Omega^1_{(A/\mathfrak p)/k}\longrightarrow0
  \tag{22.7.2.6}
\end{equation}
se sabe (20.7.20) que $g$ es sobreyectivo y que la imagen de $h$ es densa en
el núcleo de $g$ para la topología $\mathfrak m$-ádica ($\mathfrak m$ es
el ideal máximo de $A$); como $\widehat\Omega^1_{A/k}$ es un $A$-módulo de
tipo finito, todos los submódulos del $(A/\mathfrak p)$-módulo
\[
  \widehat\Omega^1_{A/k}\otimes_A(A/\mathfrak p)
  =\widehat\Omega^1_{A/k}\widehat\otimes_A(A/\mathfrak p)
\]
son cerrados para la topología $\mathfrak m$-ádica ($0_{\mathrm I}$, 7.3.5),
por lo que la sucesión (22.7.2.6) es exacta. Tensorizando por $L$, se obtiene la sucesión
exacta
\[
  (\mathfrak p/\mathfrak p^2)\otimes_AL\xrightarrow{i}
  \widehat\Omega^1_{A/k}\otimes_AL\longrightarrow
  \widehat\Omega^1_{(A/\mathfrak p)/k}\otimes_{A/\mathfrak p}L
  \longrightarrow0.
\]

Ahora bien, la hipótesis de separabilidad hecha sobre $L$ implica, en virtud de (21.9.5),
que $\widehat\Omega^1_{(A/\mathfrak p)/k}\otimes_{A/\mathfrak p}L$ tiene
rango $s$ sobre $L$; como $\widehat\Omega^1_{A/k}\otimes_AL$ tiene rango $r$
sobre $L$ (21.9.3), se tiene $\operatorname{rg}_L(\operatorname{Im}i)=r-s$. Pero como
$A/\mathfrak p$ es finita sobre una subálgebra isomorfa a
$k[[T_1,\ldots,T_s]]$, se tiene $\dim(A/\mathfrak p)=s$ ((17.1.4) y (16.1.5));
por tanto,
\[
  r-s=\dim(A)-\dim(A/\mathfrak p)=\dim(A_{\mathfrak p})
\]

\oldpage[0\textsubscript{IV-1}]{211}

por (16.5.11). Ahora bien, como $A_{\mathfrak p}$ es regular (17.3.2),
$\dim(A_{\mathfrak p})$ es igual al rango sobre $L$ de
$\mathfrak pA_{\mathfrak p}/\mathfrak p^2A_{\mathfrak p}$ (17.1.1), y por tanto al
rango sobre $L$ de $(\mathfrak p/\mathfrak p^2)\otimes_AL$, lo que acaba de
probar que $i$ es inyectivo.
\end{proof}

\begin{theorem}[22.7.3]
\label{0.22.7.3-es}
Sean $k$ un cuerpo, $A$ un anillo local noetheriano completo de cuerpo residual
$K$. Se supone que:
\begin{enumerate}
  \item[$1^\circ$] $[k:k^p]<+\infty$ (donde $p$ es el exponente característico
  de $k$);
  \item[$2^\circ$] $K$ es una extensión finita de una extensión separable
  $K_0$ de $k$;
  \item[$3^\circ$] $A$ está provisto de una estructura de $K_0$-álgebra formalmente
  lisa (para la topología preádica).
\end{enumerate}
Sean $\mathfrak q$ un ideal de $A$, $B=A/\mathfrak q$, $\mathfrak p$ un
ideal primo de $A$ que contiene $\mathfrak q$. Las condiciones siguientes son
equivalentes:
\begin{enumerate}
  \item[a)] $B_{\mathfrak p}$ es una $k$-álgebra formalmente lisa (para la
  topología $\mathfrak p$-preádica).
  \item[b)] Existen $k$-derivaciones $D_i$ de $A$ en sí mismo
  $(1\leq i\leq m)$ y elementos $f_i$ $(1\leq i\leq m)$ de
  $\mathfrak q$, tales que las imágenes de los $f_i$ en
  $\mathfrak qA_{\mathfrak p}$ engendran este ideal de $A_{\mathfrak p}$, y
  se tiene $\det(D_i(f_j))\notin\mathfrak p$.
  \item[b')] Existe una subextensión $k'$ de $K_0$, que contiene $K_0^p$,
  tal que $[K_0:k']<+\infty$, $k'$-derivaciones $D_i$ de $A$ en
  sí mismo $(1\leq i\leq m)$ y elementos $f_i$ de $\mathfrak q$
  $(1\leq i\leq m)$, tales que las imágenes de los $f_i$ en
  $\mathfrak qA_{\mathfrak p}$ engendran este ideal de $A_{\mathfrak p}$, y
  se tiene $\det(D_i(f_j))\notin\mathfrak p$.
\end{enumerate}
\end{theorem}

\begin{proof}
Distingamos dos casos, según que $p=1$ o $p>1$.

\emph{A)} $k$ \emph{es de característica} $0$. Como $K$ es entonces una
extensión separable de $K_0$, resulta de (19.6.4) que $A$ es
$K_0$-isomorfo a un anillo de series formales $K[[T_1,\ldots,T_r]]$ provisto de
su estructura usual de $K$-álgebra. Pero entonces, teniendo en cuenta (19.8.8, (ii))
aplicado a $A/\mathfrak p$, se ve que las condiciones generales de (22.7.2)
se verifican sustituyendo allí $k$ por $K$. Además, en virtud de (19.6.4), es
equivalente decir que $B_{\mathfrak p}$ es una $k$-álgebra
formalmente lisa o una $K$-álgebra formalmente lisa, pues ambas condiciones
equivalen a que $B_{\mathfrak p}$ sea un anillo regular. Se pueden
aplicar, pues, las conclusiones de (22.7.2), y es inmediato que esto
prueba la equivalencia de a), b) y b') (con $k'=K_0$ en b')).

\emph{B)} $k$ \emph{es de característica} $p>0$. Como $A$ es una
$K_0$-álgebra formalmente lisa y $K_0$ es separable sobre $k$,
resulta de (19.3.5, (ii)) y (19.6.1) que $A$ es una $k$-álgebra
formalmente lisa; en virtud de (22.5.9), $A_{\mathfrak p}$ es también una
$k$-álgebra formalmente lisa para la topología $\mathfrak p$-preádica. Resulta,
pues, de nuevo de (22.6.2, (ii)) que la condición a) equivale a que
el homomorfismo $j_0$ definido en (22.7.2.1) (donde $L=k(\mathfrak p)$) sea
inyectivo. Por tanto, (19.1.12, c)) muestra ya que b) implica a); como, por
otra parte, b') implica trivialmente b), todo se reduce a demostrar que la hipótesis
de que $j_0$ es inyectivo implica b').

Designemos primero por $k'$ una subextensión cualquiera de $K_0$ tal que
$[K_0:k']<+\infty$. Observemos que $\widehat\Omega^1_{A/k'}$ es un $A$-módulo
libre de tipo finito. Por el razonamiento de (22.7.2), basta demostrar que
(para una elección conveniente de $k'$), el homomorfismo compuesto
\begin{equation}
\label{0.22.7.3.1-es}
  j':(\mathfrak q/\mathfrak q^2)\otimes_AL
  \xrightarrow{j_0}\Omega^1_{A/k}\otimes_AL
  \longrightarrow\widehat\Omega^1_{A/k'}\otimes_AL
  \tag{22.7.3.1}
\end{equation}

\oldpage[0\textsubscript{IV-1}]{212}

es inyectivo. Consideremos de nuevo el diagrama conmutativo
\begin{equation}
\label{0.22.7.3.2-es}
\xymatrix{
  (\mathfrak q/\mathfrak q^2)\otimes_AL
    \ar[r]^{j_0} \ar[d]_\alpha &
  \Omega^1_{A/k}\otimes_AL \ar[r] \ar@{=}[d] &
  \widehat\Omega^1_{A/k'}\otimes_AL \ar@{=}[d] \\
  (\mathfrak p/\mathfrak p^2)\otimes_AL
    \ar[r]_{i_0} &
  \Omega^1_{A/k}\otimes_AL \ar[r] &
  \widehat\Omega^1_{A/k'}\otimes_AL
}
\tag{22.7.3.2}
\end{equation}

Como $j_0$ es inyectivo, se tiene $\operatorname{Ker}(i_0\circ\alpha)=0$. Pero como
$A_{\mathfrak p}$ es un anillo regular y la hipótesis a) implica que
$B_{\mathfrak p}=A_{\mathfrak p}/\mathfrak qA_{\mathfrak p}$ es también un
anillo regular (19.6.5), el lema (22.7.2.4) muestra que $\alpha$ es inyectivo,
de donde $\alpha^{-1}(\operatorname{Ker} i_0)=0$. Si $i'$ es el compuesto de los homomorfismos de la
segunda línea de (22.7.3.2), se tiene asimismo
$\operatorname{Ker}(j')=\alpha^{-1}(\operatorname{Ker}(i'))$, y todo se reduce, pues, a demostrar que, para una
elección conveniente de $k'$, se tiene
\begin{equation}
\label{0.22.7.3.3-es}
  \operatorname{Ker}(i')=\operatorname{Ker}(i).
  \tag{22.7.3.3}
\end{equation}

Pero el mismo razonamiento que en (22.7.2) muestra que se tiene una sucesión exacta
\[
  (\mathfrak p/\mathfrak p^2)\otimes_AL\xrightarrow{i'}
  \widehat\Omega^1_{A/k'}\otimes_AL\longrightarrow
  \widehat\Omega^1_{(A/\mathfrak p)/k'}\otimes_{A/\mathfrak p}L
  \longrightarrow0.
\]
Ahora bien, las hipótesis hechas permiten aplicar (21.9.8), por lo que existe una
subextensión $k'$ de $K_0$ que contiene $K_0^p$, tal que
$[K_0:k']<+\infty$, y para la cual se tiene
\[
  \operatorname{rg}_L\bigl(\widehat\Omega^1_{(A/\mathfrak p)/k'}
  \otimes_{A/\mathfrak p}L\bigr)
  =\dim(A/\mathfrak p)+\operatorname{rg}_L\Upsilon_{L/k}
  +\operatorname{rg}_{K_0}\Omega^1_{K_0/k'}.
\]
Pero se tiene
\[
  \dim(A/\mathfrak p)=\dim(A)-\dim(A_{\mathfrak p})
\]
por (16.5.11),
\[
  \dim(A_{\mathfrak p})
  =\operatorname{rg}_L\bigl((\mathfrak p/\mathfrak p^2)\otimes_AL\bigr)
\]
puesto que $A_{\mathfrak p}$ es regular (17.1.1), y finalmente
\[
  \dim(A)+\operatorname{rg}_{K_0}\Omega^1_{K_0/k'}
  =\operatorname{rg}_L\bigl(\widehat\Omega^1_{A/k'}\otimes_AL\bigr)
\]
en virtud de (21.9.2); se obtiene, pues,
\begin{equation}
\label{0.22.7.3.4-es}
  \operatorname{rg}_L(\operatorname{Ker}(i'))=\operatorname{rg}_L\Upsilon_{L/k}.
  \tag{22.7.3.4}
\end{equation}

Por otra parte, siendo el cuerpo residual de $A_{\mathfrak p}$ formalmente liso
sobre el cuerpo primo de $k$, se tiene una sucesión exacta (20.6.22.1)
\[
  \Upsilon_{L/k}\xrightarrow{\chi_{A_{\mathfrak p}/k}}
  (\mathfrak pA_{\mathfrak p})/(\mathfrak pA_{\mathfrak p})^2
  \xrightarrow{i}\Omega^1_{A_{\mathfrak p}/k}
  \otimes_{A_{\mathfrak p}}L,
\]
y como $A_{\mathfrak p}$ es una $k$-álgebra formalmente lisa para la
topología $\mathfrak p$-preádica, resulta de (22.2.2) que
$\chi_{A_{\mathfrak p}/k}$ es un homomorfismo inyectivo, por lo que $\operatorname{Ker}(i)$ es
isomorfo a $\Upsilon_{L/k}$ y, teniendo en cuenta (22.7.3.4) y que
$\operatorname{Ker}(i)\subset\operatorname{Ker}(i')$, esto acaba de probar (22.7.3.3) y, por consiguiente, el
teorema.
\end{proof}

\oldpage[0\textsubscript{IV-1}]{213}

\begin{remark}[22.7.4]
\label{0.22.7.4-es}
De hecho se ha demostrado (utilizando (21.9.6)) que la condición b') sigue siendo
equivalente a las demás condiciones de (22.7.3) cuando se obliga a $k'$ a
ser uno de los cuerpos de una familia filtrante decreciente $(k_\alpha)$ de
subcuerpos de $K_0$ que contienen $K_0^p$, tales que
$[K_0:k_\alpha]<+\infty$ para todo $\alpha$ y que
$\bigcap_\alpha k_\alpha=K_0^p$.
\end{remark}

\begin{corollary}[22.7.5]
\label{0.22.7.5-es}
Bajo las hipótesis generales de (22.7.3), el conjunto de los ideales primos
$\mathfrak n\in\operatorname{Spec}(B)$ tales que $B_{\mathfrak n}$ sea una
$k$-álgebra formalmente lisa (para la topología $\mathfrak n$-preádica)
es abierto en $\operatorname{Spec}(B)$.
\end{corollary}

\begin{proof}
Con las notaciones de (22.7.3), basta ver que, si $B_{\mathfrak p}$
es formalmente lisa sobre $k$, ocurre lo mismo con $B_{\mathfrak p'}$, para
todo $\mathfrak p'\in V(\mathfrak q)$ suficientemente próximo a $\mathfrak p$ en
$\operatorname{Spec}(A)$. Ahora bien, si las $k$-derivaciones $D_i$ y los elementos
$f_i$ de $\mathfrak q$ verifican el criterio b) de (22.7.3), se tiene también
$f=\det(D_if_j)\notin\mathfrak p'$ para $\mathfrak p'\in D(f)$; por otra
parte, existe un $g\notin\mathfrak p$ tal que las imágenes de los $f_i$ en
$\mathfrak qA_{\mathfrak p'}$ engendran $\mathfrak qA_{\mathfrak p'}$ para
$\mathfrak p'\in D(g)$ (Bourbaki, \emph{Alg. comm.}, chap.~II, §~5,
n\textsuperscript{o}~1, prop.~2), lo que acaba de probar el corolario.
\end{proof}

\begin{corollary}[22.7.6]
\label{0.22.7.6-es}
\textnormal{(Nagata).} Sea $B$ un anillo local noetheriano completo que contiene
un cuerpo. Entonces el conjunto de los $\mathfrak n\in\operatorname{Spec}(B)$ tales
que $B_{\mathfrak n}$ sea regular es abierto en $\operatorname{Spec}(B)$.
\end{corollary}

\begin{proof}
En efecto, $B$ es un álgebra sobre un cuerpo primo $k$, por tanto (19.8.8, (ii))
de la forma $A/\mathfrak q$, donde
$A=K_0[[T_1,\ldots,T_n]]$ es un anillo de series formales y $K_0$ una
extensión de $k$; por otra parte, decir que $B_{\mathfrak n}$ es regular
equivale a decir que es una $k$-álgebra formalmente lisa para la topología
$\mathfrak n$-preádica (19.6.4). Se verifican, pues, todas las condiciones de aplicación de
(22.7.5).
\end{proof}

\begin{remarks}[22.7.7]
\label{0.22.7.7-es}
\begin{enumerate}
  \item[(i)] Veremos más adelante, con Nagata, utilizando (22.7.6), que
  se puede extender la conclusión de (22.7.6) al caso de un anillo local
  noetheriano completo cualquiera (IV, 6.12.7).

  \item[(ii)] La conclusión del teorema (22.7.3) ya no es necesariamente
  exacta cuando no se supone que $[k:k^p]$ sea finito. Tomemos, por
  ejemplo, $k=\mathbf F_p(X_n)_{n\geqslant0}$,
  $A=k[[T,U]]$, $\mathfrak q$ igual al ideal principal $Af$, donde
  \[
    f=U^p-\sum_{n=0}^{\infty}X_nT^{np}~;
  \]
  $A$ es un anillo factorial en el cual $f$ es un elemento irreducible, pues es
  irreducible en el anillo de polinomios $k((T))[U]$, y por tanto también en
  el anillo de series formales $k((T))[[U]]$ (Bourbaki,
  \emph{Alg. comm.}, chap.~VII, §~3), y \emph{a fortiori} en $A$.
  Tomemos $\mathfrak p=\mathfrak q$, de modo que $B_{\mathfrak p}$ es el
  cuerpo designado por $E$ en (21.9.7), que es separable sobre $k$
  (\emph{loc. cit.}), luego una $k$-álgebra formalmente lisa. Pero se tiene aquí,
  con las notaciones de (22.7.3), $k=K_0=K$, y las $k$-derivaciones de $A$
  en sí mismo son las combinaciones lineales de $\partial/\partial T$ y
  $\partial/\partial U$; como se tiene
  $\partial f/\partial T=\partial f/\partial U=0$, los criterios b) y b') de
  (22.7.3) no se verifican.
\end{enumerate}
\end{remarks}
